% main.tex -- GENERATED, do not edit.  Regenerate with tools/make-arxiv.sh.
%
% Assembled from, in order:
%   article/preamble.tex        (typography; edit there)
%   article/frontmatter.tex     (title block and abstract; edit there)
%   the Typst source, converted by tools/typ2tex.sh
%
% Build:  xelatex main.tex  (twice, for cross-references and citations)
\documentclass[11pt,a4paper]{scrartcl}
% preamble.tex -- ALL typography for the manuscript lives here.
%
% Changelog (reverse-chronological)
%   2026-09-03  JTS, minor edits for readability
%   2026-09-01  arXiv build.  (a) Dropped "number within=section" from the
%               theorem counter: Typst numbers ALL theorem-like environments
%               globally and continuously (Theorem 1, Remark 2, Definition 3,
%               ...), and the section-scoped LaTeX numbering silently produced a
%               different sequence from the one the article's own prose, the
%               Lean sources and the project notes all cite.
%               (b) Added tikz (for the one hand-drawn figure, which pandoc
%               cannot translate) and a heading-less thebibliography, because
%               the body already carries its own numbered "References" section
%               heading converted from the Typst.
%   2026-08-30  Created, porting the Typst theorion/cosmos.fancy look to LaTeX.
%
% The point of this file, per Jan: submitting to a journal should be a ONE-FILE
% swap.  Nothing about fonts, colours, box geometry or theorem styling appears
% in the manuscript body -- the body uses only \begin{theorem}...\end{theorem}
% and friends, which every journal class already provides.
%
% To retarget:  set \documentclass to the journal's, delete the tcolorbox block
% below, and the body compiles unchanged with the class's own theorem styles.

% ---- class -------------------------------------------------------------
% KOMA (scrartcl) rather than amsart: better control of headings and margins,
% and it is not a journal class, so nothing here pretends to be a submission
% format.  An amsart variant is kept at preamble-amsart.tex.
\usepackage[a4paper,margin=2.6cm]{geometry}
\usepackage{amsmath,amssymb,amsthm,mathtools}
\usepackage{fontspec}                 % xelatex
\usepackage{microtype}
\usepackage[hidelinks]{hyperref}
\usepackage[capitalise,noabbrev]{cleveref}
\usepackage{booktabs}
% pandoc emits longtable for every table, plus these helpers; without them the
% build dies with "Environment longtable undefined" and a cascade of
% "Misplaced \noalign".
\usepackage{longtable,array,calc}
\usepackage{graphicx}
% Converted prose is dense with unbreakable inline maths, which leaves TeX no
% legal break in some lines; pandoc's own default template sets this for the
% same reason.  Without it, a handful of lines run ~25pt into the margin.
\setlength{\emergencystretch}{3em}
\providecommand{\tightlist}{\setlength{\itemsep}{0pt}\setlength{\parskip}{0pt}}
\providecommand{\pandocbounded}[1]{#1}
\usepackage[most]{tcolorbox}
\tcbuselibrary{theorems}
\usepackage{varwidth}                 % for tcolorbox's "varwidth boxed title"
\usepackage{tikz}
\usetikzlibrary{arrows.meta,positioning,shapes.geometric}

% ---- bibliography without its own heading ------------------------------
% The converted body already carries "\section{References}" (Typst's "=
% References", numbered 10.).  Stock thebibliography would add a SECOND,
% unnumbered "References" heading right under it.  This is the standard
% article.cls environment with the \section*{\refname} line removed.
\makeatletter
\renewenvironment{thebibliography}[1]
  {\list{\@biblabel{\@arabic\c@enumiv}}%
        {\settowidth\labelwidth{\@biblabel{#1}}%
         \leftmargin\labelwidth
         \advance\leftmargin\labelsep
         \@openbib@code
         \usecounter{enumiv}%
         \let\p@enumiv\@empty
         \renewcommand\theenumiv{\@arabic\c@enumiv}}%
   \sloppy
   \clubpenalty4000
   \@clubpenalty \clubpenalty
   \widowpenalty4000%
   \sfcode`\.\@m}
  {\def\@noitemerr
    {\@latex@warning{Empty `thebibliography' environment}}%
   \endlist}
\providecommand{\@openbib@code}{}
\makeatother

% ---- the coloured boxed theorem style ---------------------------------
% NOTE the empty {} prefix on every \newtcbtheorem below.  It does NOT make the
% labels bare -- tcolorbox prepends a colon unconditionally, so
% \begin{definition}{t}{def-y1y2} defines ":def-y1y2" (checked in main.aux, not
% assumed).  An empty prefix only shortens "prp:foo" to ":foo".  Cross-references
% are therefore written by tools/markers2env.py as
% \hyperref[:foo]{Definition~\ref*{:foo}}, which also supplies the environment
% name -- every box here shares one counter, so cleveref could only ever call
% all of them "Theorem".  See fix_references' docstring for the full argument.
% Reproduces theorion's cosmos.fancy: a tinted body, a saturated title bar,
% rounded corners, a coloured left rule.
% An attached boxed title is a single unbreakable line by default, so a long
% Typst title just runs off the paper: Remark 27's ("... component-by-component
% is what Proposition 26 answers") was cut at "is what P", silently -- the rest
% is in the page stream, outside the trim.  varwidth lets it wrap.
\tcbset{varwidth boxed title*=-1cm}

\definecolor{clrThm}{HTML}{2F6FAF}
\definecolor{clrDef}{HTML}{1F7A5A}
\definecolor{clrRem}{HTML}{8A6D1F}
\definecolor{clrExa}{HTML}{6B4C9A}

% NO "number within=section": Typst numbers every theorem-like environment from
% one global, continuous counter (Theorem 1, Remark 2, Definition 3, Example 4,
% Problem 5, Theorem 6, ...), and those numbers are cited by name elsewhere.
\newtcbtheorem{theorem}{Theorem}{%
  enhanced, breakable, colback=clrThm!4, colframe=clrThm, coltitle=white,
  fonttitle=\bfseries, arc=2pt, boxrule=0.6pt, left=6pt, right=6pt,
  attach boxed title to top left={xshift=6pt,yshift=-2mm},
  boxed title style={colback=clrThm, arc=1pt}}{}

% proposition/lemma share the theorem colour and counter
\newtcbtheorem[use counter from=theorem]{proposition}{Proposition}{%
  enhanced, breakable, colback=clrThm!4, colframe=clrThm, coltitle=white,
  fonttitle=\bfseries, arc=2pt, boxrule=0.6pt, left=6pt, right=6pt,
  attach boxed title to top left={xshift=6pt,yshift=-2mm},
  boxed title style={colback=clrThm, arc=1pt}}{}
\newtcbtheorem[use counter from=theorem]{lemma}{Lemma}{%
  enhanced, breakable, colback=clrThm!4, colframe=clrThm, coltitle=white,
  fonttitle=\bfseries, arc=2pt, boxrule=0.6pt, left=6pt, right=6pt,
  attach boxed title to top left={xshift=6pt,yshift=-2mm},
  boxed title style={colback=clrThm, arc=1pt}}{}
\newtcbtheorem[use counter from=theorem]{definition}{Definition}{%
  enhanced, breakable, colback=clrDef!4, colframe=clrDef, coltitle=white,
  fonttitle=\bfseries, arc=2pt, boxrule=0.6pt, left=6pt, right=6pt,
  attach boxed title to top left={xshift=6pt,yshift=-2mm},
  boxed title style={colback=clrDef, arc=1pt}}{}
\newtcbtheorem[use counter from=theorem]{remark}{Remark}{%
  enhanced, breakable, colback=clrRem!4, colframe=clrRem, coltitle=white,
  fonttitle=\bfseries, arc=2pt, boxrule=0.6pt, left=6pt, right=6pt,
  attach boxed title to top left={xshift=6pt,yshift=-2mm},
  boxed title style={colback=clrRem, arc=1pt}}{}
\newtcbtheorem[use counter from=theorem]{example}{Example}{%
  enhanced, breakable, colback=clrExa!4, colframe=clrExa, coltitle=white,
  fonttitle=\bfseries, arc=2pt, boxrule=0.6pt, left=6pt, right=6pt,
  attach boxed title to top left={xshift=6pt,yshift=-2mm},
  boxed title style={colback=clrExa, arc=1pt}}{}

\definecolor{clrPrb}{HTML}{A0522D}
\newtcbtheorem[use counter from=theorem]{problem}{Problem}{%
  enhanced, breakable, colback=clrPrb!4, colframe=clrPrb, coltitle=white,
  fonttitle=\bfseries, arc=2pt, boxrule=0.6pt, left=6pt, right=6pt,
  attach boxed title to top left={xshift=6pt,yshift=-2mm},
  boxed title style={colback=clrPrb, arc=1pt}}{}

\definecolor{clrQue}{HTML}{B03060}
% Question is the ONE environment deliberately NOT on the theorem counter, so it
% prints "Question 1" rather than taking a number out of the main sequence.  The
% Typst side does the same thing (make-frame with no counter: argument), and the
% two must agree or the routes disagree on every number after it.
%
% Reason: the Question was broken out of Problem 5 in 2026-09.  Joining the shared
% counter would have renumbered Theorem 6 -> 7, Example 7 -> 8, ... Remark 45 -> 46,
% and those numbers are cited by name in LEAN/KCommuting/*.lean, LEAN/README.md,
% tools/markers2env.py, the comment above, and Jan's review annotations.
% LEAN/README.md already documents an earlier renumbering that misled a reader.
\newtcbtheorem{question}{Question}{%
  enhanced, breakable, colback=clrQue!4, colframe=clrQue, coltitle=white,
  fonttitle=\bfseries, arc=2pt, boxrule=0.6pt, left=6pt, right=6pt,
  attach boxed title to top left={xshift=6pt,yshift=-2mm},
  boxed title style={colback=clrQue, arc=1pt}}{}

% proof stays plain -- a boxed proof reads badly at this density
\renewcommand{\qedsymbol}{$\square$}

% ---- cross-references into the boxes -----------------------------------
% \thmref{Definition}{def-y1y2}  ->  "Definition 14", hyperlinked.
% The supplement is passed in because every box above shares ONE counter
% (tcb@cnt@theorem), so cleveref has no way to tell a Lemma from a Remark; and
% the ":" is tcolorbox's own, prepended to the third argument regardless of the
% \newtcbtheorem prefix.
% Deliberately NO square brackets in the call: an earlier version emitted
% \hyperref[:foo]{...} directly, and amsthm builds its proof heading as
% \item[...#1...], whose optional-argument scanner ends at the first "]" -- the
% one inside \hyperref's own [..].  204 errors, all from that.
\newcommand{\thmref}[2]{\hyperref[:#2]{#1~\ref*{:#2}}}

% ---- displays that are wider than the measure --------------------------
% Typst shrinks an over-wide display to fit (Example 38's 3x9 Jacobian, split
% into a 3x6 block, is typeset visibly smaller there).  LaTeX instead lets it
% run off the paper: 144pt of the matrix was outside the trim, with only an
% Overfull \hbox in the log to say so.  \fitdisplay measures first and scales
% ONLY when the content really is too wide, so ordinary displays are untouched.
\makeatletter
\newsavebox{\jts@fitbox}
\newcommand{\fitdisplay}[1]{%
  \sbox{\jts@fitbox}{\ensuremath{\displaystyle #1}}%
  \ifdim\wd\jts@fitbox>\linewidth
    \resizebox{\linewidth}{!}{\usebox{\jts@fitbox}}%
  \else
    \usebox{\jts@fitbox}%
  \fi}
\makeatother

% ---- notation ----------------------------------------------------------
\newcommand{\FF}{\mathbb{F}}
\newcommand{\QQ}{\mathbb{Q}}
\newcommand{\RR}{\mathbb{R}}
\newcommand{\AAA}{\mathbb{A}}
\newcommand{\PP}{\mathbb{P}}
\newcommand{\vek}[1]{\overline{#1}}

% Filenames like code/verify\_generator\_independence.sage occur in running
% prose and have no legal break point, so TeX pushes 148pt of one into the
% margin.  Allow a break after each underscore; \allowbreak sets nothing, so
% the filename is unchanged, it just becomes breakable.
\let\jtsOldUnderscore\_
\renewcommand{\_}{\jtsOldUnderscore\allowbreak}
\newcommand{\enc}{\operatorname{enc}}
\newcommand{\red}{\operatorname{red}}
\DeclareMathOperator{\tr}{tr}
\DeclareMathOperator{\pd}{pd}

\begin{document}

% frontmatter.tex -- title block and abstract, transcribed from the Typst source.
%
% Changelog (reverse-chronological)
%   2026-09-01  Created.
%
% tools/typ2tex.sh strips everything before the first "= " heading (the Typst
% preamble imports theorion, which pandoc's evaluator cannot run), and the title
% block, author block, ORCID, affiliation, e-mail, abstract and title footnote
% all live in that stripped region.  They are transcribed here BY HAND from
% article/ring-of-k-commuting-2x2-matrices.typ, which makes this file a second
% copy that can drift from the canonical source -- it did, within a day of being
% written.  Edit the .typ first, then mirror the change here, and rely on
% make-arxiv.sh's internal-matter check to catch what mirroring misses.  The
% Acknowledgements paragraph is NOT here: it sits after the bibliography in the
% Typst and survives the conversion on its own.

\definecolor{orcidgreen}{HTML}{A6CE39}
\newcommand{\orcidid}[1]{%
  \href{https://orcid.org/#1}{%
    \textcolor{orcidgreen}{\fontsize{8}{9}\selectfont iD}%
    \,\fontsize{8}{9}\selectfont #1}}

% NB: the title is NOT wrapped in a minipage.  A \footnote inside a minipage is
% typeset at the foot of the minipage, not the page -- the title footnote would
% land immediately under the title.
\begingroup
\centering
% "plus 1fil" on BOTH sides, not bare lengths: \centering works by setting
% \leftskip/\rightskip to 0pt plus 1fil, so overwriting them with rigid lengths
% narrows the measure but also destroys the centring -- the title's last line
% came out stretched across the full measure.
\leftskip=0.05\textwidth plus 1fil \rightskip=0.05\textwidth plus 1fil
    \fontsize{15}{18}\selectfont\bfseries
    The coordinate ring of the $k$-fold iterated commutator locus for
    $2 \times 2$ matrices%
    \par
\endgroup

\begingroup
\centering
\vspace{0.8em}
{\fontsize{11}{13}\selectfont Jan Snellman\hspace{0.25em}%
  \orcidid{0009-0002-6676-5068}\par}

\vspace{0.35em}

{\fontsize{9}{11}\selectfont Matematiska Institutionen, Linköpings
  Universitet, 581 83 Linköping, Sweden\par}

\vspace{0.15em}

{\fontsize{9}{11}\selectfont
  \href{mailto:jan.snellman@liu.se}{jan.snellman@liu.se}\par}
\endgroup

\vspace{1.0em}

\begin{center}
\begin{minipage}{\dimexpr\textwidth-3.2cm\relax}
  \fontsize{9.2}{11.5}\selectfont
  \begin{center}\textbf{Abstract}\end{center}
  \vspace{0.25em}
  For $2 \times 2$ matrices $A_1,\dots,A_k$, write $[A_1,\dots,A_k]$ for the
  left-normed iterated commutator $[\dots[[A_1,A_2],A_3],\dots,A_k]$, and $I_k$
  for the ideal, in the $3k$-variable reduced-coordinate polynomial ring $R_k$,
  cutting out its vanishing locus. We prove, for every $k \geq 2$ over any field
  of characteristic $\neq 2$, and as four independently-established results
  rather than one bundled claim: $I_k$ has codimension 2; $I_k$ has exactly 3
  minimal generators; $R_k / I_k$ is Cohen-Macaulay; and $I_k$ is radical. The
  last of these, together with an explicit component count resting on a
  non-containment argument, assembles into the \textbf{Primary Decomposition
  Theorem}: $I_k = P_2 \cap \cdots \cap P_k$ is an irredundant primary
  decomposition into exactly $k-1$ primes, following an explicit recursive
  block-involvement pattern. The proof identifies $I_k$ as the ideal of
  $2 \times 2$ minors of an explicit $2 \times 3$ matrix (a determinantal ideal,
  not merely determinantal-\textit{looking}), and invokes classical
  determinantal-ideal theory (Bruns-Vetter) and an explicit rank-2 Jacobian
  witness on every component (Serre's criterion) for the algebraic and
  radicality halves respectively.

  \vspace{0.9em}
  {\fontsize{8.6}{10.5}\selectfont
   \textbf{Keywords.} Iterated commutator; determinantal ideal; commuting
   variety; Cohen-Macaulay ring; primary decomposition; radical ideal; cross
   product; $\mathfrak{sl}_2$.\par}

  \vspace{0.3em}
  {\fontsize{8.6}{10.5}\selectfont
   \textbf{2020 Mathematics Subject Classification.} Primary 13C40, 14M12;
   Secondary 13H10, 15A27, 13P10.\par}
\end{minipage}
\end{center}

\vspace{0.8em}

\section{Introduction}

\(2 \times 2\) matrices \(A,B\) have commutator \(\lbrack A,B\rbrack = AB - BA\), and the \textbf{iterated} left-normed commutator of \(k\) matrices is \(\left\lbrack A_{1},\ldots,A_{k} \right\rbrack \coloneqq  \left\lbrack \ldots\left\lbrack \left\lbrack A_{1},A_{2} \right\rbrack,A_{3} \right\rbrack,\ldots,A_{k} \right\rbrack\), matching the standard lower-central-series convention for Lie algebras. The object arose for the author twice: in the fiber-counting theory of word maps \(\pi:F_{n} \rightarrow G\), \(F_{n}\) the free group on \(n\) letters, after a close reading of Tambour \cite{Tambour2000} -- and in the algebraic-complexity question of how many scalar multiplications are needed to \textbf{compute} \(\left\lbrack A_{1},\ldots,A_{k} \right\rbrack\), which remains open for \(k \geq 3\). The one fact needed from that story is purely algebraic: an explicit bilinear formula for the ordinary \(2 \times 2\) commutator, chained across \(k\) matrices, gives an explicit polynomial parametrization of \(\left\lbrack A_{1},\ldots,A_{k} \right\rbrack\)'s reduced coordinates (\thmref{Definition}{def-wk} below), and the subject of this article is the \textbf{vanishing locus} of that parametrization, as an algebraic variety, together with its coordinate ring. At \(k = 2\) this locus \textbf{is} the classical commuting variety; for \(k \geq 3\) the vanishing of the iterated bracket is weaker than pairwise commutation, so the locus studied here properly contains the pairwise-commuting one.

The classical object -- \(C_{n} \coloneqq  \left\{ (A,B) \in M_{n} \times M_{n}:AB = BA \right\}\) for \(n \times n\) matrices -- has been studied since the 1950s. Its irreducibility was proven independently by Motzkin and Taussky \cite{MotzkinTaussky1955} and by Gerstenhaber \cite{Gerstenhaber1961}, the latter also establishing the bound \(\dim_{\mathbb{F}}\mathbb{F}\lbrack A,B\rbrack \leq n\) for the subalgebra generated by two commuting matrices. Whether the coordinate ring \(\mathbb{F}\lbrack A,B\rbrack/I\), \(I\) generated by the entries of \(AB - BA\), is \textbf{reduced} and \textbf{Cohen-Macaulay} for every \(n\) is a much harder question, attributed to Michael Artin and Melvin Hochster \cite{ArtinHochsterConjecture} and stated precisely as ``Conjecture 1.1'' in Majidi-Zolbanin and Snapp \cite{MajidiZolbaninSnapp}, while Chau \cite{Chau2024} records a dated (1984) reducedness conjecture of Hochster about the same loci. Hreinsdóttir \cite{Hreinsdottir2005} attacks the conjecture case by case at \(n = 3,4\); Knutson \cite{Knutson2003} studies the minimal primes of closely related schemes. At the time of writing, Cohen-Macaulayness is settled for \(n \leq 4\) (by hand at \(n \leq 2\), by machine computation at \(n = 3,4\) \cite{MajidiZolbaninSnapp} \cite{Hreinsdottir2005}) and open for every \(n \geq 5\); reducedness is settled for \(n \leq 3\) and open for \(n \geq 4\) \cite{Chau2024}.

At \(n = 2\) -- this article's \(k = 2\) base case -- the question is comparatively easy, and is answered again, as a special case, by \thmref{Theorem}{thm-cm} below: \(I_{2}\) (\thmref{Definition}{def-wk}) is literally the ideal of the \(n = 2\) commuting variety, expressed in reduced (trace-removed) coordinates. The genuinely new object is the \textbf{iterated} generalization to \(k \geq 3\): not the classical commuting variety, but a structurally analogous determinantal ideal, built by chaining the same bilinear cross-product formula that governs the \(n = 2\) case -- and for it, the same two questions the Artin-Hochster conjecture asks, Cohen-Macaulayness and reducedness, are answered completely and unconditionally for every \(k \geq 2\) (\thmref{Theorem}{thm-primarydecomp}). In that sense the results below are a fully worked structural analogue of the still-open classical question, in a setting simple enough (\(n = 2\), but \(k\) unbounded) to admit a complete answer.

\section{Setup: the signed encoding and the defining ideal}\label{sec-setup}

Throughout, \(\mathbb{F}\) is a field of characteristic \(\neq 2\), and all matrices, vectors and polynomial rings are over \(\mathbb{F}\) unless said otherwise; the displayed examples below are computed over \(\mathbb{Q}\), but nothing in the definitions is special to it. One typographic convention is used throughout: an \textbf{overline} denotes a \textbf{vector} -- an element of \(\mathbb{F}^{3}\) (or, for \(R_{k}\) itself, of the free module \(R_{k}^{3}\)) -- while the same letter, plain, with a two-index comma subscript, denotes one of that vector's individual scalar \textbf{components}. So \({\overline{v}}_{i} \in \mathbb{F}^{3}\) is block \(i\)'s whole reduced-coordinate vector, with components \(v_{i,1},v_{i,2},v_{i,3} \in \mathbb{F}\); likewise \({\overline{w}}_{j}\) versus \(w_{j,1},w_{j,2},w_{j,3}\).

\textbf{Step 1: discard the scalar part.} Write \(A = \begin{pmatrix}
a_{11} & a_{12} \\
a_{21} & a_{22}
\end{pmatrix}\) as the sum of its scalar part \(\left( \text{tr}(A)/2 \right)I\) and its traceless remainder \begin{equation}\fitdisplay{A_{0} \coloneqq  A - \left( \text{tr}(A)/2 \right)I = \begin{pmatrix}
d_{a}/2 & a_{12} \\
a_{21} & - d_{a}/2
\end{pmatrix},\quad d_{a} \coloneqq  a_{11} - a_{22}.}\end{equation} Scalar matrices are central, so they never contribute to a commutator: \begin{equation}\fitdisplay{\lbrack cI,B\rbrack = cIB - BcI = cB - cB = 0}\end{equation} for every \(B\) and every scalar \(c\), hence \(\lbrack A,B\rbrack = \left\lbrack A_{0},B_{0} \right\rbrack\) identically -- no information relevant to \textbf{any} bracket is lost by discarding the scalar part, so from here on only the traceless remainder matters.

\textbf{Step 2: encode the traceless remainder.} A traceless matrix \(\begin{pmatrix}
x & y \\
z & - x
\end{pmatrix}\) has three free entries. Rather than reading them off directly, take \begin{equation}\fitdisplay{\text{ enc}(A) \coloneqq  \left( a_{11} - a_{22},a_{12} + a_{21},a_{12} - a_{21} \right) \in \mathbb{F}^{3},}\end{equation} the traceless \((1,1)\)-entry doubled. On a traceless matrix written as above this reads \begin{equation}\fitdisplay{\text{ enc }\begin{pmatrix}
x & y \\
z & - x
\end{pmatrix} = (2x,y + z,y - z),}\end{equation} and both forms are worth having in view: the first is what one applies to a general \(A\), the second is what one reads when a point of \(\mathbb{A}^{3}\) has to be produced or checked by hand. (On the traceless part \(a_{22} = - a_{11}\), so \(d_{a} = a_{11} - a_{22} = 2a_{11}\): the quantity \(d_{a}\) of Step 1 is literally twice that entry.) The remaining two coordinates are \(A\)'s two off-diagonal entries replaced by their sum and difference. No division is needed to \textbf{compute} \(\text{enc}\) -- only its inverse (recovering \(a_{12},a_{21}\) individually as \(\frac{v_{2} + v_{3}}{2}\) and \(\frac{v_{2} - v_{3}}{2}\)) needs \(2\) invertible.

\textbf{Step 3: form the commutator.} Given traceless \(A_{0} = \begin{pmatrix}
x & y \\
z & - x
\end{pmatrix}\) and \(B_{0} = \begin{pmatrix}
u & v \\
w & - u
\end{pmatrix}\), the ordinary matrix commutator \(\left\lbrack A_{0},B_{0} \right\rbrack = A_{0}B_{0} - B_{0}A_{0}\) is again traceless automatically, since \(\text{tr}(PQ) = \text{tr}(QP)\) for any \(P,Q\) and therefore \(\text{tr}\left( \lbrack P,Q\rbrack \right) = 0\) always, so it is again encodable by \(\text{enc}\).

\textbf{Step 4: the encoded bracket is the cross product, up to one fixed sign twist.}

For the following theorem, the article's standing hypothesis \(\text{char }\mathbb{F} \neq 2\) is deliberately not used: the identity holds over every commutative ring.

\begin{theorem}{the basic identity}{thm-basic-identity} For traceless \(2 \times 2\) matrices \(A_{0},B_{0}\) over \(\mathbb{F}\), with no hypothesis on the characteristic, \begin{equation}\fitdisplay{\text{ enc}\left( \left\lbrack A_{0},B_{0} \right\rbrack \right) = {S\left( \text{enc}\left( A_{0} \right) \times \text{ enc}\left( B_{0} \right) \right)},\quad S \coloneqq  \begin{pmatrix}
 - 1 & 0 & 0 \\
0 & - 1 & 0 \\
0 & 0 & 1
\end{pmatrix}.}\end{equation}

\end{theorem}

\begin{proof}[{Proof of \thmref{Theorem}{thm-basic-identity}}] The proof is a direct computation, carried out symbolically over generic entries \(a_{11},\ldots,b_{22}\) rather than checked at sampled points, in code/verify\_lorentzian\_encoding\_identity.sage. Both sides, expanded in the 8 free entries, agree term by term; they are displayed in full in \cref{app-identity}.

Both sides are integer polynomials in those entries, so the identity descends to every commutative ring. It was also checked directly over \(\mathbb{Z}\) and, separately, over \(\mathbb{F}_{2}\) (where it still holds), in code/referee\_round4\_setup\_determinantal\_check.sage.

\end{proof}

\begin{definition}{reduced iterated commutator, signed encoding}{def-wk} \begin{equation}\fitdisplay{{\overline{w}}_{2} \coloneqq  {\overline{v}}_{1} \times {\overline{v}}_{2},\quad{\overline{w}}_{j} \coloneqq  {S\left( {\overline{w}}_{j - 1} \right)} \times {\overline{v}}_{j}\quad(j = 3,\ldots,k).}\end{equation} Let \(R_{k} \coloneqq  \mathbb{F}\left\lbrack {\overline{v}}_{1},\ldots,{\overline{v}}_{k} \right\rbrack\) (\({\overline{v}}_{i} \in \mathbb{F}^{3}\), \(3k\) variables) and \(I_{k} \subset R_{k}\) the ideal generated by the three components of \({\overline{w}}_{k}\) -- the ideal cutting out \(\left\lbrack A_{1},\ldots,A_{k} \right\rbrack = 0\) in reduced coordinates. Two deliberate understatements are worth flagging. First, \(I_{k}\) is at this point only the \textbf{generated} ideal; that it is also the full vanishing ideal of that locus is \thmref{Theorem}{thm-radical} below, and is not to be assumed before then. Second, the encoded bracket is strictly \(S\left( {\overline{w}}_{k} \right)\), not \({\overline{w}}_{k}\) -- but \(S\) merely negates two of three coordinates, so the two triples generate the same ideal, and it is cleaner to carry the untwisted one.

\end{definition}

To spell out exactly what this means: ``\(R_{k} \coloneqq  \mathbb{F}\left\lbrack {\overline{v}}_{1},\ldots,{\overline{v}}_{k} \right\rbrack\)'' is shorthand for \textbf{the polynomial ring in the \(3k\) scalar variables \(v_{i,j}\) (\(i = 1,\ldots,k\), \(j = 1,2,3\))}, grouped into \(k\) blocks of 3 -- \textbf{not} a ring with \(k\) vector-valued generators (a polynomial ring's generators are always scalars). At \(k = 2\), written out in full: \begin{equation}\fitdisplay{R_{2} = \mathbb{F}\left\lbrack v_{1,1},v_{1,2},v_{1,3},v_{2,1},v_{2,2},v_{2,3} \right\rbrack,}\end{equation} six ordinary scalar indeterminates, with \({\overline{v}}_{1} = \left( v_{1,1},v_{1,2},v_{1,3} \right)\) and \({\overline{v}}_{2} = \left( v_{2,1},v_{2,2},v_{2,3} \right)\) merely a notational grouping of them into two blocks of three -- convenient because \({\overline{w}}_{2} = {\overline{v}}_{1} \times {\overline{v}}_{2}\) is most naturally written as a cross product of two block-vectors, not as some formula in six unrelated scalars.

Each \({\overline{w}}_{j}\) is a triple of polynomials, i.e. \({\overline{w}}_{j} \in R_{j}^{3}\) (the free rank-3 module over \(R_{j}\)) -- equivalently, the polynomial map \({\overline{w}}_{j}:\mathbb{A}^{3j} \rightarrow \mathbb{A}^{3}\) \textbf{given by substitution}, evaluating each of \({\overline{w}}_{j}\)'s three polynomial components at a point of \(\mathbb{A}^{3j}\); and \(I_{k}\) (\thmref{Definition}{def-wk}) is, equivalently, the ideal generated by the three coordinate functions of \({\overline{w}}_{k}\).

\begin{example}{\({\overline{w}}_{k}\) and \(I_{k}\), explicitly, for \(k = 2,3,4\)}{ex-wk-signed} Everything below was computed directly (code/compute\_signed\_encoding\_examples.sage) rather than derived by hand. At \(k = 2\), the three components of \({\overline{w}}_{2} = {\overline{v}}_{1} \times {\overline{v}}_{2}\): \begin{equation}\fitdisplay{v_{1,2}v_{2,3} - v_{1,3}v_{2,2},\quad v_{1,3}v_{2,1} - v_{1,1}v_{2,3},\quad v_{1,1}v_{2,2} - v_{1,2}v_{2,1}.}\end{equation} At \(k = 3\), the three components of \({\overline{w}}_{3} = {S\left( {\overline{w}}_{2} \right)} \times {\overline{v}}_{3}\) -- with \textbf{every coefficient equal to \(\pm 1\)}, no factor of 2 anywhere: \begin{equation}\fitdisplay{v_{1,2}v_{2,1}v_{3,2} - v_{1,1}v_{2,2}v_{3,2} - v_{1,3}v_{2,1}v_{3,3} + v_{1,1}v_{2,3}v_{3,3},}\end{equation} \begin{equation}\fitdisplay{- v_{1,2}v_{2,1}v_{3,1} + v_{1,1}v_{2,2}v_{3,1} - v_{1,3}v_{2,2}v_{3,3} + v_{1,2}v_{2,3}v_{3,3},}\end{equation} \begin{equation}\fitdisplay{v_{1,3}v_{2,1}v_{3,1} - v_{1,1}v_{2,3}v_{3,1} + v_{1,3}v_{2,2}v_{3,2} - v_{1,2}v_{2,3}v_{3,2},}\end{equation} and \(I_{3} \coloneqq  \left( w_{3,1},w_{3,2},w_{3,3} \right)\). At \(k = 4\), \({\overline{w}}_{4} = {S\left( {\overline{w}}_{3} \right)} \times {\overline{v}}_{4}\): again every coefficient is \(\pm 1\); the first component of \({\overline{w}}_{4}\) is \begin{equation}\fitdisplay{- v_{1,3}v_{2,1}v_{3,1}v_{4,2} + v_{1,1}v_{2,3}v_{3,1}v_{4,2} - v_{1,3}v_{2,2}v_{3,2}v_{4,2} + v_{1,2}v_{2,3}v_{3,2}v_{4,2}}\end{equation} \begin{equation}\fitdisplay{+ v_{1,2}v_{2,1}v_{3,1}v_{4,3} - v_{1,1}v_{2,2}v_{3,1}v_{4,3} + v_{1,3}v_{2,2}v_{3,3}v_{4,3} - v_{1,2}v_{2,3}v_{3,3}v_{4,3},}\end{equation} the other two of the same shape (omitted for space; all three components of \({\overline{w}}_{4}\), and \(I_{4} \coloneqq  \left( w_{4,1},w_{4,2},w_{4,3} \right)\), in code/compute\_signed\_encoding\_examples.sage's output). The same data at \(k = 6\), where the generators reach 32 terms apiece and the \(\pm 1\) pattern is still unbroken, is in \cref{app-i6}.

\end{example}

The displayed examples above were computed over \(\mathbb{Q}\); every result below holds over the arbitrary field \(\mathbb{F}\) of characteristic \(\neq 2\) fixed at the start of this section.

\section{The determinantal structure and the generator count}\label{sec-gens}

\begin{theorem}{determinantal structure}{thm-determinantal} For \(k \geq 3\), \(I_{k}\) is the ideal of \(2 \times 2\) minors of the \(2 \times 3\) matrix \begin{equation}\fitdisplay{\begin{pmatrix}
{S\left( {\overline{w}}_{k - 1} \right)} \\
{\overline{v}}_{k}
\end{pmatrix},}\end{equation} i.e. \(V\left( I_{k} \right)\) is the preimage, under \begin{equation}\fitdisplay{\left( {\overline{v}}_{1},\ldots,{\overline{v}}_{k} \right) \mapsto \begin{pmatrix}
{S\left( {\overline{w}}_{k - 1}\left( {\overline{v}}_{1},\ldots,{\overline{v}}_{k - 1} \right) \right)} \\
{\overline{v}}_{k}
\end{pmatrix},}\end{equation} of the classical rank-\(\leq 1\) determinantal variety of \(2 \times 3\) matrices. (At \(k = 2\) the same statement holds with \(S\left( {\overline{w}}_{k - 1} \right)\) replaced by \({\overline{v}}_{1}\) itself, matching \thmref{Definition}{def-wk}'s base case \({\overline{w}}_{2} = {\overline{v}}_{1} \times {\overline{v}}_{2}\) directly.)

\end{theorem}

\begin{proof}[{Proof of \thmref{Theorem}{thm-determinantal}}] The claim is immediate from two facts taken together. The three components of a cross product \(a \times b\) are, up to sign, the \(2 \times 2\) minors of the \(2 \times 3\) matrix with rows \(a\) and \(b\); and \thmref{Definition}{def-wk}'s recursion builds \({\overline{w}}_{k}\) from exactly such a cross product. The statement was also checked independently at \(k = 2,\ldots,6\) in code/verify\_determinantal\_structure\_all\_k.sage, in two separate ways. First, \(I_{k}\) and the ideal of the displayed minors are shown equal by reducing each ideal's generators to zero against the other's Gröbner basis, in \textbf{both} directions, rather than by containment one way. Second -- and this is what makes ``independent'' mean something here -- the recursion itself is checked against honest \(2 \times 2\) matrix multiplication: the \(A_{i}\) are reconstructed from their \(\text{enc}\)-coordinates, the iterated commutator is formed by repeated matrix products, and the result encoded and compared against \(S\left( {\overline{w}}_{k} \right)\). The two routes share no code. The underlying polynomial identity \(\text{enc}\left( \lbrack A,B\rbrack \right) = {S\left( \text{enc}(A) \times \text{ enc}(B) \right)}\) of \thmref{Theorem}{thm-basic-identity} makes this a two-line induction valid for \textbf{every} \(k\), not merely a pattern checked at small \(k\) (code/verify\_lorentzian\_encoding\_identity.sage).

\end{proof}

This theorem is the engine behind the geometric results of \cref{sec-geom}: \(R_{k}/I_{k}\) is not merely \textbf{determinantal-looking}, it is literally an iterated extension of the coordinate ring of a classical determinantal variety.

\begin{example}{the \(2 \times 3\) matrix of \thmref{Theorem}{thm-determinantal}, for \(k = 2,3,4\)}{auto1} At \(k = 2\): \begin{equation}\fitdisplay{\begin{pmatrix}
{\overline{v}}_{1} \\
{\overline{v}}_{2}
\end{pmatrix} = \begin{pmatrix}
v_{1,1} & v_{1,2} & v_{1,3} \\
v_{2,1} & v_{2,2} & v_{2,3}
\end{pmatrix}.}\end{equation} At \(k = 3\), the first row is \({S\left( {\overline{w}}_{2} \right)} = \left( - w_{2,1}, - w_{2,2},w_{2,3} \right)\) where \({\overline{w}}_{2} = {\overline{v}}_{1} \times {\overline{v}}_{2}\) is as in the previous example, giving \begin{equation}\fitdisplay{\begin{pmatrix}
{S\left( {\overline{w}}_{2} \right)} \\
{\overline{v}}_{3}
\end{pmatrix} = \begin{pmatrix}
v_{1,3}v_{2,2} - v_{1,2}v_{2,3} & - v_{1,3}v_{2,1} + v_{1,1}v_{2,3} & - v_{1,2}v_{2,1} + v_{1,1}v_{2,2} \\
v_{3,1} & v_{3,2} & v_{3,3}
\end{pmatrix}.}\end{equation} At \(k = 4\) the first row is \(S\left( {\overline{w}}_{3} \right)\), whose three entries are cubic with four terms each. Naming them \(A\), \(B\), \(C\) keeps the matrix legible: \begin{equation}\fitdisplay{\begin{pmatrix}
A & B & C \\
v_{4,1} & v_{4,2} & v_{4,3}
\end{pmatrix},}\end{equation} where \begin{equation}\fitdisplay{\begin{aligned}
A & = - v_{1,2}v_{2,1}v_{3,2} + v_{1,1}v_{2,2}v_{3,2} + v_{1,3}v_{2,1}v_{3,3} - v_{1,1}v_{2,3}v_{3,3}, \\
B & = v_{1,2}v_{2,1}v_{3,1} - v_{1,1}v_{2,2}v_{3,1} + v_{1,3}v_{2,2}v_{3,3} - v_{1,2}v_{2,3}v_{3,3}, \\
C & = v_{1,3}v_{2,1}v_{3,1} - v_{1,1}v_{2,3}v_{3,1} + v_{1,3}v_{2,2}v_{3,2} - v_{1,2}v_{2,3}v_{3,2}.
\end{aligned}}\end{equation} The row degree grows by one at each step (\(1,2,3,\ldots\)), while the second row is always the fresh, free vector \({\overline{v}}_{k}\) -- exactly the ``one fixed complicated row, one free row'' shape \thmref{Proposition}{prop-detideal} needs.

\end{example}

\begin{lemma}{non-degeneracy, with an explicit witness}{lem-nondeg} For every \(j \geq 2\), the polynomial map \({\overline{w}}_{j}:\mathbb{A}^{3j} \rightarrow \mathbb{A}^{3}\) is not identically zero. Concretely, the point \begin{equation}\fitdisplay{p^{(j)} \coloneqq  \left( e_{1},e_{2},\underset{j - 2}{\underbrace{e_{1},\ldots,e_{1}}} \right) \in \mathbb{A}^{3j}}\end{equation} satisfies \({\overline{w}}_{j}\left( p^{(j)} \right) \neq 0\); explicitly \({\overline{w}}_{j}\left( p^{(j)} \right) = e_{3}\) for \(j\) even and \(e_{2}\) for \(j\) odd. Hence \(U_{j} \coloneqq  \left\{ p \in \mathbb{A}^{3j}:{\overline{w}}_{j}(p) \neq 0 \right\}\) is a nonempty Zariski-open subset of \(\mathbb{A}^{3j}\).

\end{lemma}

\begin{proof}[{Proof of \thmref{Lemma}{lem-nondeg}}] The induction is on \(j\), and it carries more than the statement ``some point survives'': it carries a \textbf{named} point, because \thmref{Lemma}{lem-witness-new} later needs one rather than a mere existence claim.

\textbf{Base case \(j = 2\).} By \thmref{Definition}{def-wk}, \({\overline{w}}_{2} = {\overline{v}}_{1} \times {\overline{v}}_{2}\), so \({\overline{w}}_{2}\left( e_{1},e_{2} \right) = e_{1} \times e_{2} = e_{3} \neq 0\). Thus \(p^{(2)} = \left( e_{1},e_{2} \right)\) works.

\textbf{Induction step.} Let \(j \geq 3\) and suppose \(p \in \mathbb{A}^{3(j - 1)}\) satisfies \(c \coloneqq  {\overline{w}}_{j - 1}(p) \neq 0\). The step rests on three observations, taken in order:

\begin{enumerate}
\item
  \({S(c)} \neq 0\), since \(\det S = 1\) makes \(S\) invertible over every field, so it kills no nonzero vector.
\item
  Therefore \(\mathbb{F} \cdot {S(c)}\) is a \textbf{line} in \(\mathbb{F}^{3}\). The three vectors \(e_{1},e_{2},e_{3}\) span \(\mathbb{F}^{3}\), and a line contains at most one of them; so at least two indices \(i\) have \(e_{i} \notin \mathbb{F} \cdot {S(c)}\). Fix any such \(i\).
\item
  A cross product vanishes exactly on linearly dependent pairs (\(\ker\lbrack u\rbrack_{\times} = \mathbb{F}u\) for \(u \neq 0\), the computation carried out in the proof of \thmref{Lemma}{lem-witness-inherited}), so \({S(c)} \times e_{i} \neq 0\).
\end{enumerate}

By \thmref{Definition}{def-wk}, \({\overline{w}}_{j}\left( p,e_{i} \right) = {S\left( {\overline{w}}_{j - 1}(p) \right)} \times e_{i} = {S(c)} \times e_{i} \neq 0\). (The twist falls on \(c\) rather than on the result because that is where \thmref{Definition}{def-wk}'s recursion puts it.) So \({\overline{w}}_{j}\) is not identically zero, and appending \(e_{i}\) to a witness at level \(j - 1\) gives one at level \(j\).

\textbf{The witness in closed form.} Running the step from \(p^{(2)} = \left( e_{1},e_{2} \right)\) and taking the least admissible \(i\) each time, the choice is \(i = 1\) at \textbf{every} level: \({\overline{w}}_{2} = e_{3}\) and \({S\left( e_{3} \right)} = e_{3}\), which is not proportional to \(e_{1}\), giving \({\overline{w}}_{3} = e_{3} \times e_{1} = e_{2}\); then \({S\left( e_{2} \right)} = - e_{2}\), again not proportional to \(e_{1}\), giving \({\overline{w}}_{4} = \left( - e_{2} \right) \times e_{1} = e_{3}\); and the pair \(\left( e_{3},e_{2} \right)\) now repeats. Hence the closed form in the statement, and every coordinate of both the witness and its image lies in \(\{ 0, \pm 1\}\). This has been verified for \(j = 2,\ldots,8\), by two independent evaluations per level, in code/verify\_nondegeneracy\_witness.sage, and the results are tabulated in \cref{app-witness}.

\textbf{Openness.} \(U_{j}\) is the complement of the Zariski-closed set \(V\left( w_{j,1},w_{j,2},w_{j,3} \right)\), hence open; it is nonempty because \(p^{(j)}\) lies in it.

\end{proof}

The generator count of \(I_{k}\) needs nothing beyond the definition and non-degeneracy -- not codimension, not the determinantal-structure theorem above, not any external citation:

\begin{lemma}{independent same-degree generators are minimal}{lem-mingens-generic} If a homogeneous ideal in a graded ring is generated by finitely many forms of the same degree that are linearly independent over the base field, that generating set is minimal.

\end{lemma}

\begin{proof}[{Proof of \thmref{Lemma}{lem-mingens-generic}}] A polynomial-coefficient combination reproducing one of the generators from the others, while staying in the same degree, can only use degree-0 (constant) coefficients -- so a redundant generator would violate linear independence directly.

\end{proof}

\begin{lemma}{the three generators of \(I_{k}\) are independent}{lem-mingens-indep} The three components of \({\overline{w}}_{k}\) are linearly independent over \(\mathbb{F}\), for every \(k \geq 2\).

\end{lemma}

\begin{proof}[{Proof of \thmref{Lemma}{lem-mingens-indep}}] By induction on \(k\), and the induction hypothesis must be \textbf{independence} itself; see \thmref{Remark}{rem-nondeg-insufficient} for why the weaker non-degeneracy of \thmref{Lemma}{lem-nondeg} does not suffice.

\textbf{Base case \(k = 2\).} The three components of \({\overline{w}}_{2} = {\overline{v}}_{1} \times {\overline{v}}_{2}\) are the \(2 \times 2\) minors of the generic \(2 \times 3\) matrix with rows \({\overline{v}}_{1},{\overline{v}}_{2}\), displayed in \thmref{Example}{ex-wk-signed}. Each contains a monomial the other two do not: \(v_{1,2}v_{2,3}\) occurs only in the first, \(v_{1,3}v_{2,1}\) only in the second, and \(v_{1,1}v_{2,2}\) only in the third. So no nontrivial constant relation is possible.

\textbf{Induction step.} Suppose the components of \({\overline{w}}_{k - 1}\) are linearly independent over \(\mathbb{F}\), and let \(c \in \mathbb{F}^{3}\) satisfy \(c \cdot {\overline{w}}_{k} = 0\). Writing \(a \coloneqq  {S\left( {\overline{w}}_{k - 1} \right)}\), the scalar triple product identity gives \begin{equation}\fitdisplay{0 = c \cdot \left( a \times {\overline{v}}_{k} \right) = (c \times a) \cdot {\overline{v}}_{k}.}\end{equation} The variables \({\overline{v}}_{k}\) are a \textbf{fresh} block, occurring in neither \(a\) nor \(c\), and the right-hand side is linear in them; so each of its three coefficients vanishes separately, i.e. \begin{equation}\fitdisplay{c \times a = 0.}\end{equation}

Suppose \(c \neq 0\). Then \(c \times a = 0\) forces \(a = f \cdot c\) for a single polynomial \(f\): over the fraction field \(c \times a = 0\) says \(a\) lies in the line spanned by \(c\), and \(f = a_{i}/c_{i}\) is a polynomial for any index with \(c_{i} \neq 0\). But then \(a\)'s three components are pairwise proportional, hence linearly \textbf{dependent} over \(\mathbb{F}\) (any \(d\) with \(d \cdot c = 0\) gives \(d \cdot a = (d \cdot c)f = 0\), and such a nonzero \(d\) exists since \(c\) spans only a line in \(\mathbb{F}^{3}\)). Since \(S\) is invertible, the components of \({\overline{w}}_{k - 1} = S^{- 1}(a)\) are then linearly dependent too, contradicting the induction hypothesis. Hence \(c = 0\).

Corroborated computationally: the coefficient matrix has rank 3 at \(k = 2,\ldots,6\) (code/verify\_generator\_independence.sage); the full table is in \cref{app-independence}.

This lemma is also \textbf{formally verified}: it is the statement w\_linearIndependent of the Lean 4 / mathlib development accompanying this article, proved there for every \(k \geq 2\) over an arbitrary field rather than checked at small \(k\). \thmref{Theorem}{thm-mingens} below is formalized too, as w\_minimal\_generators. \cref{sec-methods} says what the formalization does and does not cover.

\end{proof}

\begin{remark}{why non-degeneracy is not enough here}{rem-nondeg-insufficient} It is tempting to argue instead that a constant relation would force \(S\left( {\overline{w}}_{k - 1} \right)\) to be a \textbf{constant vector}, and to rule that out by \thmref{Lemma}{lem-nondeg}. That argument is not valid. The relation forces only \(c \times {S\left( {\overline{w}}_{k - 1} \right)} = 0\), which makes the components of \(S\left( {\overline{w}}_{k - 1} \right)\) pairwise proportional, not constant. A vector such as \((f,2f,3f)\) with \(f \neq 0\) is nonvanishing, so \thmref{Lemma}{lem-nondeg} does not exclude it, yet its components are visibly dependent. Non-degeneracy is strictly weaker than what the step needs, which is why the induction above carries linear independence itself.

\end{remark}

\begin{theorem}{\(I_{k}\) has exactly 3 minimal generators}{thm-mingens} For every \(k \geq 2\): \(I_{k}\) has exactly 3 minimal generators.

\end{theorem}

\begin{proof}[{Proof of \thmref{Theorem}{thm-mingens}}] Immediate from \thmref{Lemma}{lem-mingens-generic} applied to \thmref{Lemma}{lem-mingens-indep}'s independence, since \(I_{k}\) is generated by the three (same-degree-\(k\)-homogeneous) components of \({\overline{w}}_{k}\) by definition.

\end{proof}

\section{\texorpdfstring{The geometric structure of \(V\left( I_{k} \right)\)}{The geometric structure of V\textbackslash left( I\_\{k\} \textbackslash right)}}\label{sec-geom}

\begin{lemma}{continuous images and closures preserve irreducibility}{lem-top-irred} Let \(f:X \rightarrow Y\) be a continuous map of topological spaces with \(X\) irreducible. Then \(f(X)\) (with the subspace topology) is irreducible, and hence so is \(\overline{f(X)}\).

\end{lemma}

\begin{proof}[{Proof of \thmref{Lemma}{lem-top-irred}}] \textbf{Image.} Suppose \(f(X) = C_{1} \cup C_{2}\) with \(C_{1},C_{2}\) closed in \(f(X)\). By continuity, \(f^{- 1}\left( C_{1} \right)\) and \(f^{- 1}\left( C_{2} \right)\) are closed in \(X\), and \(X = f^{- 1}\left( C_{1} \right) \cup f^{- 1}\left( C_{2} \right)\). Since \(X\) is irreducible, one of the two equals \(X\) -- say \(f^{- 1}\left( C_{1} \right) = X\) -- so \(f(X) \subseteq C_{1}\), i.e. \(C_{1} = f(X)\), which is not proper. Hence \(f(X)\) is irreducible. \textbf{Closure.} If \(A \subseteq Y\) is irreducible and \(\overline{A} = C_{1} \cup C_{2}\) with \(C_{1},C_{2}\) closed in \(Y\), then \(A = \left( A \cap C_{1} \right) \cup \left( A \cap C_{2} \right)\); irreducibility of \(A\) forces, say, \(A \subseteq C_{1}\); since \(C_{1}\) is closed, \(\overline{A} \subseteq C_{1}\), and combined with \(C_{1} \subseteq \overline{A}\), \(C_{1} = \overline{A}\), not proper. Apply both parts with \(A \coloneqq  f(X)\).

\end{proof}

This is purely topological -- continuity of \(f\) is the only hypothesis used, nothing about \(f\) being a polynomial map or morphism of varieties \cite{Youcis2014}.

\begin{definition}{the strata \(Y_{1}\), \(Y_{2}\)}{def-y1y2} For \(k \geq 3\), write \(\pi\) for the projection \textbf{forgetting \({\overline{v}}_{k}\)}, \begin{equation}\fitdisplay{\pi:\mathbb{A}^{3k} \rightarrow \mathbb{A}^{3(k - 1)},\quad\pi({\overline{v}}_{1},\ldots,{\overline{v}}_{k}) \coloneqq  \left( {\overline{v}}_{1},\ldots,{\overline{v}}_{k - 1} \right),}\end{equation} and \(\Phi\) for the map \begin{equation}\fitdisplay{\Phi:\mathbb{A}^{3(k - 1)} \times \mathbb{A}^{1} \rightarrow \mathbb{A}^{3k},\quad\Phi(p,\lambda) \coloneqq  \left( p,\lambda{S\left( {\overline{w}}_{k - 1}(p) \right)} \right).}\end{equation} Define two subsets of \(\mathbb{A}^{3k}\): \begin{equation}\fitdisplay{Y_{1} \coloneqq  V\left( I_{k - 1} \right) \times \mathbb{A}^{3},\quad Y_{2} \coloneqq  \overline{\Phi(\mathbb{A}^{3k - 2})}.}\end{equation} Here \(Y_{1}\) is the locus where the previous bracket already vanishes, with \({\overline{v}}_{k}\) completely free, while \(Y_{2}\) is the closure of the locus swept out as \({\overline{v}}_{k}\) ranges over scalar multiples of \(S\left( {\overline{w}}_{k - 1}(p) \right)\), over all \(p \in \mathbb{A}^{3(k - 1)}\). At \(k = 2\) there is no ``previous bracket'' to chain from, so neither \(Y_{1}\) nor \(\Phi\) is defined; by convention \(Y_{2} \coloneqq  V\left( I_{2} \right)\) directly there, matching \thmref{Theorem}{thm-determinantal}'s own base-case treatment.

\end{definition}

\begin{lemma}{\(Y_{2}\) is irreducible}{lem-y2-irred} For \(k \geq 3\): \(Y_{2}\) (\thmref{Definition}{def-y1y2}) is irreducible, and in fact its defining ideal (the graph ideal of \(\Phi\) with \(\lambda\) eliminated) is prime.

\end{lemma}

\begin{proof}[{Proof of \thmref{Lemma}{lem-y2-irred}}] The map \(\Phi\) is polynomial, hence continuous for the Zariski topology, since preimages of Zariski-closed sets under a polynomial map are again Zariski-closed by substitution. Its domain \(\mathbb{A}^{3k - 2}\) is irreducible, because the polynomial ring over \(\mathbb{F}\) is a domain and so its zero ideal is prime. \thmref{Lemma}{lem-top-irred} then gives \(Y_{2}\) irreducible directly. Primality of the elimination ideal follows the same way: it is the kernel of a map into a domain (\(\mathbb{F}\left\lbrack {\overline{v}}_{1},\ldots,{\overline{v}}_{k} \right\rbrack/( \cdot ) \cong \mathbb{F}\left\lbrack {\overline{v}}_{1},\ldots,{\overline{v}}_{k - 1},\lambda \right\rbrack\)), hence prime -- this holds for every \(k\) by this argument alone, additionally corroborated computationally at \(k = 4\) by computing the elimination ideal directly and confirming primality in Sage, as an independent check rather than a substitute for the argument.

\end{proof}

\begin{example}{the defining ideals of \(Y_{1}\) \textbf{and} \(Y_{2}\), for \(k = 2,3,4\)}{auto2} Writing \({\overline{v}}_{i} = \left( v_{i,1},v_{i,2},v_{i,3} \right)\) for the three coordinates of block \({\overline{v}}_{i}\): at \(k = 2\), \(Y_{2} \coloneqq  V\left( I_{2} \right)\) (the base-case convention above), so its ideal \textbf{is} \(I_{2}\), generated by \begin{equation}\fitdisplay{v_{1,2}v_{2,3} - v_{1,3}v_{2,2},\quad v_{1,3}v_{2,1} - v_{1,1}v_{2,3},\quad v_{1,1}v_{2,2} - v_{1,2}v_{2,1}.}\end{equation} (\(Y_{1}\) is not defined at \(k = 2\), \thmref{Definition}{def-y1y2}.) For \(k \geq 3\), \(Y_{1}\)'s defining ideal is simply \(I_{k - 1}\) itself, viewed inside \(R_{k}\) by leaving \({\overline{v}}_{k}\) free -- at \(k = 3\), exactly the three generators displayed above for \(I_{2}\), now inside \(R_{3}\) with \({\overline{v}}_{3}\) not appearing at all; at \(k = 4\), exactly \(I_{3}\)'s three generators (displayed below in this same example) inside \(R_{4}\) with \({\overline{v}}_{4}\) free. Unlike \(Y_{2}\), no elimination is needed to compute \(Y_{1}\)'s ideal -- it is already \(I_{k - 1}\)'s own generating set, unchanged. At \(k = 3\), \(Y_{2}\)'s elimination ideal (computed directly, not inferred) is generated by \begin{equation}\fitdisplay{v_{2,1}v_{3,1} + v_{2,2}v_{3,2} - v_{2,3}v_{3,3},\quad v_{1,1}v_{3,1} + v_{1,2}v_{3,2} - v_{1,3}v_{3,3},}\end{equation} \begin{equation}\fitdisplay{v_{1,2}v_{2,1}v_{3,2} - v_{1,1}v_{2,2}v_{3,2} - v_{1,3}v_{2,1}v_{3,3} + v_{1,1}v_{2,3}v_{3,3}}\end{equation} -- and these three repay a careful reading, because their shape is not the one the parametrization first suggests. \textbf{All three} are linear in the new block \({\overline{v}}_{3}\); what separates them is total degree, \(2,2,3\). The two quadratic ones are, up to sign, the incidence conditions \begin{equation}\fitdisplay{{\overline{v}}_{1} \cdot {S\left( {\overline{v}}_{3} \right)} = 0,\quad{\overline{v}}_{2} \cdot {S\left( {\overline{v}}_{3} \right)} = 0.}\end{equation} These are \textbf{consequences} of the proportionality \({\overline{v}}_{3} \propto {S\left( {\overline{w}}_{2} \right)}\), not its components: substituting \({\overline{v}}_{3} = \lambda{S\left( {\overline{w}}_{2} \right)}\) and using \(S^{2} = I\) turns them into \(\lambda\left( {\overline{v}}_{i} \cdot {\overline{w}}_{2} \right) = 0\), which holds because \({\overline{w}}_{2} = {\overline{v}}_{1} \times {\overline{v}}_{2}\) is orthogonal to both of its factors. The proportionality condition itself -- the vanishing of \({\overline{v}}_{3} \times {S\left( {\overline{w}}_{2} \right)}\), whose three components are, up to sign, exactly the components of \({\overline{w}}_{3}\) -- enters through a single representative: the third generator, of degree 3, is verbatim the first component of \({\overline{w}}_{3}\), and the other two components of \({\overline{w}}_{3}\) also lie in this ideal (as they must, since \(I_{3} \subseteq P_{3}\); confirmed by direct reduction) without being needed as generators. Note that no generator of this ideal could avoid \({\overline{v}}_{3}\) altogether: \(Y_{2}\) projects dominantly onto \(\mathbb{A}^{6}\), so the elimination ideal meets \(R_{2}\) in \(0\). \textbf{This ideal strictly contains \(I_{3}\)} -- \(I_{3}\) is its intersection with a second prime (the ``\(Y_{1}\)'' piece, where the \(k = 2\) bracket already vanishes); see \thmref{Example}{ex-components-k3}, where this exact ideal reappears as one of \(I_{3}\)'s two primary components. At \(k = 4\) the same pattern repeats one degree higher, the three generators now having degrees \(2,3,4\) and again all three linear in the new block \({\overline{v}}_{4}\). The first two are once more incidence conditions forced by \({\overline{w}}_{3} \perp {S\left( {\overline{w}}_{2} \right)},{\overline{v}}_{3}\), namely \({\overline{v}}_{3} \cdot {S\left( {\overline{v}}_{4} \right)} = 0\) and \({\overline{w}}_{2} \cdot {\overline{v}}_{4} = 0\); the third is the negative of \({\overline{w}}_{4}\)'s first component: \begin{equation}\fitdisplay{v_{3,1}v_{4,1} + v_{3,2}v_{4,2} - v_{3,3}v_{4,3},}\end{equation} \begin{equation}\fitdisplay{v_{1,3}v_{2,2}v_{4,1} - v_{1,2}v_{2,3}v_{4,1} - v_{1,3}v_{2,1}v_{4,2} + v_{1,1}v_{2,3}v_{4,2} + v_{1,2}v_{2,1}v_{4,3} - v_{1,1}v_{2,2}v_{4,3},}\end{equation} \begin{equation}\fitdisplay{v_{1,3}v_{2,1}v_{3,1}v_{4,2} - v_{1,1}v_{2,3}v_{3,1}v_{4,2} + v_{1,3}v_{2,2}v_{3,2}v_{4,2} - v_{1,2}v_{2,3}v_{3,2}v_{4,2}}\end{equation} \begin{equation}\fitdisplay{- v_{1,2}v_{2,1}v_{3,1}v_{4,3} + v_{1,1}v_{2,2}v_{3,1}v_{4,3} - v_{1,3}v_{2,2}v_{3,3}v_{4,3} + v_{1,2}v_{2,3}v_{3,3}v_{4,3}.}\end{equation} All computed directly (elimination of \(\lambda\) from the graph ideal of \(\Phi\), in Sage), not just asserted from the pattern.

\end{example}

\begin{lemma}{dimension of \(Y_{2}\), and where it lives}{lem-y2-dim} With \(Y_{2}\) as in \thmref{Lemma}{lem-y2-irred} and \(U_{k - 1}\) as in \thmref{Lemma}{lem-nondeg}: \(\dim Y_{2} = 3k - 2\), and \(Y_{2}\) meets \(\pi^{- 1}\left( U_{k - 1} \right)\) nonemptily, where \(\pi\) is the projection of \thmref{Definition}{def-y1y2}, forgetting \({\overline{v}}_{k}\).

\end{lemma}

\begin{proof}[{Proof of \thmref{Lemma}{lem-y2-dim}}] The projection \(\pi:\mathbb{A}^{3k} \rightarrow \mathbb{A}^{3(k - 1)}\) discards the last block, so that \(\pi({\overline{v}}_{1},\ldots,{\overline{v}}_{k}) = \left( {\overline{v}}_{1},\ldots,{\overline{v}}_{k - 1} \right)\). Hence \(\pi^{- 1}\left( U_{k - 1} \right) = U_{k - 1} \times \mathbb{A}^{3}\), the set of points whose \textbf{first} \(k - 1\) blocks lie in \(U_{k - 1}\) and whose \({\overline{v}}_{k}\) is completely unconstrained. Over \(U_{k - 1}\), \(\Phi\) is injective. Both \(p\) and \(\lambda\) can be recovered from \(\Phi(p,\lambda)\): \(p\) directly, and \(\lambda\) by dividing any nonzero coordinate of \({\overline{v}}_{k} = \lambda{S\left( {\overline{w}}_{k - 1}(p) \right)}\) by the corresponding coordinate of \(S\left( {\overline{w}}_{k - 1}(p) \right)\); this division is well-defined since \({S\left( {\overline{w}}_{k - 1}(p) \right)} \neq 0\) on \(U_{k - 1}\). Injectivity gives \(\dim Y_{2} = \dim\left( U_{k - 1} \times \mathbb{A}^{1} \right) = 3(k - 1) + 1 = 3k - 2\). Since \(U_{k - 1}\) is nonempty by \thmref{Lemma}{lem-nondeg}, \(\Phi(U_{k - 1} \times \left\{ 1 \right\})\) is a nonempty subset of \(Y_{2} \cap \pi^{- 1}\left( U_{k - 1} \right)\).

\end{proof}

\begin{theorem}{\(I_{k}\) has codimension 2}{thm-codim} For every \(k \geq 2\): \(\dim V\left( I_{k} \right) = 3k - 2\), i.e. \(\text{codim}\left( I_{k} \right) = 2\).

\end{theorem}

\begin{proof}[{Proof of \thmref{Theorem}{thm-codim}}] The proof is by induction on \(k\). \textbf{Base case \(k = 2\)}: \(V\left( I_{2} \right)\) is the rank-\(\leq 1\) locus of a \(2 \times 3\) matrix with independent-indeterminate entries. Its dimension is the standard \(4 = 3 \cdot 2 - 2\) for a generic determinantal variety of this shape: the classical codimension formula for the rank-\(\leq t\) locus of a generic \(m \times n\) matrix gives \((m - t)(n - t) = 1 \cdot 2 = 2\) here \cite{BrunsVetter1988}, and the same value falls out of viewing \(V\left( I_{2} \right)\) as the affine cone over the Segre embedding \(\mathbb{P}^{1} \times \mathbb{P}^{2} \hookrightarrow \mathbb{P}^{5}\), of dimension \(1 + 2 + 1 = 4\). The value is confirmed directly by Gröbner-basis computation as well (\thmref{Example}{ex-dim-codim}). \textbf{Inductive step}: by \thmref{Theorem}{thm-determinantal}, and since \(a \times b = 0\) if and only if \(a,b\) are linearly dependent, \(V\left( I_{k} \right) = Y_{1} \cup Y_{2}\) (\thmref{Definition}{def-y1y2}). The locus \({S\left( {\overline{w}}_{k - 1} \right)} = 0\) gives \(Y_{1}\), with \({\overline{v}}_{k}\) unconstrained; the locus \({S\left( {\overline{w}}_{k - 1} \right)} \neq 0\) with \({\overline{v}}_{k}\) proportional to it gives points \textbf{literally in \(\Phi\)'s image}, no closure needed, hence in \(Y_{2}\). For the reverse containment \(Y_{1} \cup Y_{2} \subseteq V\left( I_{k} \right)\): \({\overline{w}}_{k} = {S\left( {\overline{w}}_{k - 1} \right)} \times {\overline{v}}_{k}\) vanishes wherever \({\overline{w}}_{k - 1}\) does, so \(Y_{1} \subseteq V\left( I_{k} \right)\); and \(\Phi( \cdot ) \subseteq V\left( I_{k} \right)\) identically, since a vector crossed with its own scalar multiple vanishes, so \(Y_{2} = \overline{\Phi( \cdot )} \subseteq V\left( I_{k} \right)\), the vanishing locus being closed. By the inductive hypothesis, \(\dim V\left( I_{k - 1} \right) = 3(k - 1) - 2\), so \(\dim Y_{1} = 3(k - 1) - 2 + 3 = 3k - 2\). Also \(\dim Y_{2} = 3k - 2\), by \thmref{Lemma}{lem-y2-dim}. Hence \(\dim V\left( I_{k} \right) = \dim(Y_{1} \cup Y_{2}) = 3k - 2\). This dimension statement alone needs no containment argument; that is needed only for the sharper component-count result below.

\end{proof}

\begin{example}{dimension and codimension, for \(k = 2,3,4\)}{ex-dim-codim} Values of \(\dim R_{k}\) (the ambient dimension, \(= 3k\)), \(\dim(R_{k}/I_{k})\) (the dimension of the affine variety \(V\left( I_{k} \right)\) -- \textbf{not} ``\(\dim I_{k}\)'', an ideal is not a variety and does not itself have a dimension in this sense), and \(\text{codim}\left( I_{k} \right) \coloneqq  \dim R_{k} - \dim(R_{k}/I_{k})\):

\begin{longtable}[]{@{}cccc@{}}
\toprule\noalign{}
\endhead
\bottomrule\noalign{}
\endlastfoot
\(k\) & \(\dim R_{k}\) & \(\dim(R_{k}/I_{k})\) & \(\text{codim}\left( I_{k} \right)\) \\
2 & 6 & 4 & 2 \\
3 & 9 & 7 & 2 \\
4 & 12 & 10 & 2 \\
\end{longtable}

All three rows follow directly from \thmref{Theorem}{thm-codim}'s formula \(\dim(R_{k}/I_{k}) = 3k - 2\); also confirmed by direct Gröbner-basis dimension computation in Sage at each \(k\), as an independent check.

\end{example}

\begin{proposition}{\(Y_{2} \nsubseteq Y_{1}\)}{prop-y2-not-in-y1} With \(Y_{1},Y_{2}\) as in the proof of \thmref{Theorem}{thm-codim}: \(Y_{2} \nsubseteq Y_{1}\).

\end{proposition}

\begin{proof}[{Proof of \thmref{Proposition}{prop-y2-not-in-y1}}] Pick \(p \in U_{k - 1}\) (\thmref{Lemma}{lem-nondeg}); then \(\Phi(p,0) = (p,0) \in Y_{2}\) but \({\overline{w}}_{k - 1}(p) \neq 0\) puts \((p,0) \notin Y_{1}\).

\end{proof}

\begin{proposition}{no component of \(Y_{1}\) lies in \(Y_{2}\)}{prop-y1-not-in-y2} Let \(Z\) be an irreducible component of \(V\left( I_{k - 1} \right)\) \textbf{of dimension \(3(k - 1) - 2\)}. Then \(Z \times \mathbb{A}^{3} \nsubseteq Y_{2}\).

The dimension hypothesis is stated rather than cited. \thmref{Theorem}{thm-codim} gives the dimension of the \textbf{union} \(V\left( I_{k - 1} \right)\), which does not by itself say that every component attains it; the per-component statement is part of \thmref{Theorem}{thm-componentcount}, and this proposition is used inside that theorem's own induction, where it is available as the inductive hypothesis at level \(k - 1\).

\end{proposition}

\begin{proof}[{Proof of \thmref{Proposition}{prop-y1-not-in-y2}}] Both \(Z \times \mathbb{A}^{3}\) and \(Y_{2}\) are irreducible: \(Y_{2}\) by \thmref{Lemma}{lem-y2-irred}, and \(Z \times \mathbb{A}^{3}\) because a product of irreducible varieties is irreducible (\(Z\) is irreducible by hypothesis, and \(\mathbb{A}^{3}\) trivially is). Both have the \textbf{same} dimension \(3k - 2\): \(Z \times \mathbb{A}^{3}\) by the hypothesis on \(Z\), and \(Y_{2}\) by \thmref{Lemma}{lem-y2-dim}. Containment would therefore force equality. This is refuted: \(Z \times \mathbb{A}^{3}\) is disjoint from \(\pi^{- 1}\left( U_{k - 1} \right)\), since \({\overline{w}}_{k - 1}\) vanishes identically on \(Z \subset V\left( I_{k - 1} \right)\); but \(Y_{2}\) meets \(\pi^{- 1}\left( U_{k - 1} \right)\), by \thmref{Lemma}{lem-y2-dim}.

\end{proof}

\begin{theorem}{component count}{thm-componentcount} For every \(k \geq 2\): \(V\left( I_{k} \right)\) has exactly \(k - 1\) irreducible components, each of dimension \(3k - 2\). The component indexed \(j = 2,\ldots,k\) involves exactly the first \(j\) coordinate blocks and nothing later.

\end{theorem}

\begin{proof}[{Proof of \thmref{Theorem}{thm-componentcount}}] The proof is by induction on \(k\). \textbf{Base case \(k = 2\)}: the rank-\(\leq 1\) locus of a generic \(2 \times 3\) matrix is irreducible -- the affine cone over the Segre embedding \(\mathbb{P}^{1} \times \mathbb{P}^{2} \hookrightarrow \mathbb{P}^{5}\), equivalently the Motzkin-Taussky \cite{MotzkinTaussky1955} irreducibility of the \(2 \times 2\) commuting variety -- so exactly 1 component, matching \(k - 1 = 1\). \textbf{Step}: \(V\left( I_{k} \right) = Y_{1} \cup Y_{2}\) as in \thmref{Theorem}{thm-codim}'s proof. By the inductive hypothesis (this theorem at level \(k - 1\), which is where the per-component dimension comes from), \(V\left( I_{k - 1} \right)\) has \(k - 2\) irreducible components, each of dimension \(3(k - 1) - 2\), none contained in another; crossing with \(\mathbb{A}^{3}\) makes them \(Y_{1}\)'s components, each of dimension \(3k - 2\). That is exactly the hypothesis \thmref{Proposition}{prop-y1-not-in-y2} asks for, so it and \thmref{Proposition}{prop-y2-not-in-y1} together rule out containment between \(Y_{2}\) and any of them, in either direction. A finite union of irreducible closed sets with no containments is irredundant, giving exactly \((k - 2) + 1 = k - 1\) components.

\end{proof}

This gives \(I_{k}\)'s minimal primes and their count, but not yet that \(I_{k}\) is radical (so not yet that these are literally the components of a \textbf{primary} decomposition of \(I_{k}\) itself, as opposed to of \(\sqrt{I_{k}}\)) -- that is closed in \cref{sec-rad}, after Cohen-Macaulayness is available.

\begin{definition}{the components \(P_{j}\)}{def-primecomponents} For \(k \geq 2\) and \(j = 2,\ldots,k\), let \(\Phi_{j}\) be \thmref{Definition}{def-y1y2}'s map \(\Phi\) run at parameter \(j\) (with the base convention \(\Phi_{2}(p,\lambda) \coloneqq  (p,\lambda p)\)), and let \begin{equation}\fitdisplay{\Psi_{j}:\mathbb{A}^{3(j - 1)} \times \mathbb{A}^{1} \times \mathbb{A}^{3(k - j)} \rightarrow \mathbb{A}^{3k},\quad(p,\lambda,q) \mapsto \left( \Phi_{j}(p,\lambda),q \right)}\end{equation} be \(\Phi_{j}\) with the remaining blocks \({\overline{v}}_{j + 1},\ldots,{\overline{v}}_{k}\) carried along untouched. Since \(\Psi_{j}\) is a polynomial map it \textbf{is}, read backwards, a homomorphism of \(\mathbb{F}\)-algebras \begin{equation}\fitdisplay{\varphi_{j}:R_{k} \rightarrow \mathbb{F}\lbrack p,\lambda,q\rbrack = \mathbb{F}\left\lbrack t_{1},\ldots,t_{3k - 2} \right\rbrack,}\end{equation} sending each variable \(v_{i,c}\) to the corresponding coordinate polynomial of \(\Psi_{j}\). Define \begin{equation}\fitdisplay{P_{j} \coloneqq  \ker\varphi_{j} \subset R_{k}.}\end{equation} Two consequences come for free. First, each \(P_{j}\) is \textbf{prime}, being the kernel of a homomorphism into a domain -- for every \(j\) including \(j = 2\), with no case distinction and no appeal to \thmref{Lemma}{lem-y2-irred}, which is stated only for \(k \geq 3\). Second, \(P_{j}\) does not depend on which \(k \geq j\) it is viewed inside of: enlarging \(k\) appends free variables to source and target alike.

\textbf{Why the kernel rather than the vanishing ideal.} It is tempting to define \(P_{j}\) as the vanishing ideal of the \(j\)-th component of \(V\left( I_{k} \right)\), and over \(\overline{\mathbb{F}}\) the two agree -- \thmref{Proposition}{prop-rational} identifies that component as the closure of \(\Psi_{j}\)'s image. But over a \textbf{finite} field ``vanishing ideal'' is the wrong notion outright: \(V\left( I_{k} \right)\left( \mathbb{F} \right)\) is then a finite set of points, and \(x^{q} - x\) vanishes on all of \(\mathbb{A}^{1}\left( \mathbb{F}_{q} \right)\) without being the zero polynomial. The kernel is the right object over every field at once, is well defined with no convention attached, and is what makes the base-change statement in \thmref{Remark}{rem-base-field} an identity rather than a coincidence.

\end{definition}

\begin{example}{a suitable point of \(Y_{2}\), concretely, at \(k = 3\)}{auto3} It is worth having a single explicit point, and not only the defining ideal of \thmref{Example}{ex-components-k3} below. Take \({\overline{v}}_{1} = (1,0,0)\) and \({\overline{v}}_{2} = (0,1,0)\), so that \({\overline{v}}_{1} \in U_{2}\) by \thmref{Lemma}{lem-nondeg}, giving \({\overline{w}}_{2} = {\overline{v}}_{1} \times {\overline{v}}_{2} = (0,0,1)\) and \({S\left( {\overline{w}}_{2} \right)} = (0,0,1)\). Taking \(\lambda = 1\) in \thmref{Definition}{def-y1y2}'s \(\Phi\), i.e. \({\overline{v}}_{3} \coloneqq  {S\left( {\overline{w}}_{2} \right)} = (0,0,1)\), gives the point \begin{equation}\fitdisplay{\left( {\overline{v}}_{1},{\overline{v}}_{2},{\overline{v}}_{3} \right) = \left( (1,0,0),(0,1,0),(0,0,1) \right) \in Y_{2} \subset \mathbb{A}^{9.}}\end{equation} One checks this directly: \({\overline{w}}_{3} = {S\left( {\overline{w}}_{2} \right)} \times {\overline{v}}_{3} = (0,0,1) \times (0,0,1) = 0\), a vector crossed with itself. So the point lies in \(V\left( I_{3} \right)\), as any point of the form \(\Phi(p,\lambda)\) must (\thmref{Definition}{def-y1y2}). The point lies on the \textbf{new} component \(P_{3}\), not on the inherited \(P_{2}\) -- \({\overline{v}}_{1}\) and \({\overline{v}}_{2}\) are independent here, so the \(k = 2\) bracket does not already vanish. It is \textbf{not} a rank-2 witness in \thmref{Lemma}{lem-witness-new}'s sense (that witness uses \(\lambda = 0\), not \(\lambda = 1\)) -- it illustrates \(Y_{2}\) itself, not the radicality argument's own witness points.

\end{example}

\begin{proposition}{rational parametrization of each component}{prop-rational} For \(k \geq 2\) and \(j = 2,\ldots,k\), the component of \(V\left( I_{k} \right)\) indexed \(j\) (\thmref{Theorem}{thm-componentcount}, over \(\overline{\mathbb{F}}\) per \thmref{Remark}{rem-base-field}) is exactly the closure of the image of \thmref{Definition}{def-primecomponents}' map \(\Psi_{j}\); equivalently, \(P_{j} = \ker\varphi_{j}\) is its vanishing ideal. In particular that component is a \textbf{rational} variety. The parametrization is moreover \textbf{birational} onto the component: \(\Psi_{j}\) is injective over \(U_{j - 1} \times \mathbb{A}^{1} \times \mathbb{A}^{3(k - j)}\) -- where, at \(j = 2\), the symbol \(U_{1}\) is to be read as \(\{{\overline{v}}_{1} \neq 0\}\), since \thmref{Lemma}{lem-nondeg} defines \(U_{j}\) only for \(j \geq 2\) and \({\overline{w}}_{1}\) does not exist.

\end{proposition}

\begin{proof}[{Proof of \thmref{Proposition}{prop-rational}}] At parameter \(j = k\) the map \(\Psi_{k}\) is \(\Phi\) itself, whose image closure is \(Y_{2}\) (\thmref{Definition}{def-y1y2}), and \thmref{Theorem}{thm-componentcount}'s induction identifies \(Y_{2}\) with the component indexed \(k\). For \(j < k\) the same construction run at parameter \(j\) gives that level's \(Y_{2}\), and appending the free blocks \({\overline{v}}_{j + 1},\ldots,{\overline{v}}_{k}\) crosses it with \(\mathbb{A}^{3(k - j)}\) -- which is exactly \thmref{Theorem}{thm-componentcount}'s description of the inherited component indexed \(j\). That the vanishing ideal of the closure of a polynomial map's image is the kernel of the corresponding algebra homomorphism is the standard identification, over \(\overline{\mathbb{F}}\).

For birationality, \thmref{Lemma}{lem-y2-dim}'s injectivity argument (over \(U_{j - 1}\), via \thmref{Lemma}{lem-nondeg} at level \(j - 1\)) applies verbatim at parameter \(j\), and appending the free blocks \(q\) unchanged does not affect injectivity.

\end{proof}

\begin{remark}{a reducible variety is not itself ``rational'' -- component-by-component is what \thmref{Proposition}{prop-rational} answers}{auto4} \(Y_{1} \coloneqq  V\left( I_{k - 1} \right) \times \mathbb{A}^{3}\) (\thmref{Definition}{def-y1y2}) is \textbf{not} irreducible for \(k \geq 4\) -- it is a union of \(k - 2\) of the components above (\thmref{Theorem}{thm-componentcount}'s inductive step) -- so a single rational parametrization \textbf{of \(Y_{1}\) as one object} is not meaningful in the usual sense (rationality is a property of irreducible varieties). What \textbf{is} meaningful, and what \thmref{Proposition}{prop-rational} supplies, is a rational parametrization of each of \(Y_{1}\)'s irreducible pieces individually -- exactly the components inherited from \(V\left( I_{k - 1} \right)\).

\end{remark}

\begin{example}{the components, worked out at \(k = 3\)}{ex-components-k3} Uniformly by \(k\), since the components' \textbf{count} and \textbf{block pattern} are already known (\thmref{Theorem}{thm-componentcount}, \thmref{Definition}{def-primecomponents}) and only their explicit generators change:

\begin{itemize}
\item
  \textbf{At \(k = 2\)}: \(I_{2}\) is already prime (\thmref{Theorem}{thm-componentcount}'s base case), so \(V\left( I_{2} \right)\) has exactly \(k - 1 = 1\) component -- itself.
\item
  \textbf{At \(k = 3\)}: exactly \(k - 1 = 2\) components, both computed directly (Gröbner bases, not inferred from the pattern alone), matching \(I_{3} = P_{2} \cap P_{3}\):

  \begin{itemize}
  \item
    \(P_{2}\), the ``\(Y_{1}\)'' piece (block pattern \(\{ 1,2\}\), matching that the length-2 sub-bracket already vanishes with \({\overline{v}}_{3}\) free) -- generated by exactly \(I_{2}\)'s three generators from the first example, now inside \(R_{3}\) with \({\overline{v}}_{3}\) simply not appearing: \begin{equation}\fitdisplay{v_{1,2}v_{2,3} - v_{1,3}v_{2,2},\quad v_{1,3}v_{2,1} - v_{1,1}v_{2,3},\quad v_{1,1}v_{2,2} - v_{1,2}v_{2,1}.}\end{equation}
  \item
    \(P_{3}\), the ``\(Y_{2}\)'' piece (block pattern \(\{ 1,2,3\}\)) -- exactly the ideal already computed in the \(Y_{2}\)-defining-ideal example above: \begin{equation}\fitdisplay{v_{2,1}v_{3,1} + v_{2,2}v_{3,2} - v_{2,3}v_{3,3},\quad v_{1,1}v_{3,1} + v_{1,2}v_{3,2} - v_{1,3}v_{3,3},}\end{equation} \begin{equation}\fitdisplay{v_{1,2}v_{2,1}v_{3,2} - v_{1,1}v_{2,2}v_{3,2} - v_{1,3}v_{2,1}v_{3,3} + v_{1,1}v_{2,3}v_{3,3}.}\end{equation}
  \end{itemize}
\item
  \textbf{At \(k = 4\)}: exactly \(k - 1 = 3\) components, block patterns \(\{ 1,2\}\), \(\{ 1,2,3\}\), \(\{ 1,2,3,4\}\):

  \begin{itemize}
  \item
    The first two, \(P_{2}\) and \(P_{3}\), are inherited \textbf{verbatim} from \(k = 3\) above, with \({\overline{v}}_{4}\) simply left free -- not recomputed, per \thmref{Definition}{def-primecomponents}' own ``does not depend on which \(k\)'' clause.
  \item
    The third, \(P_{4}\), is the new ``\(Y_{2}\)'' stratum -- exactly the \(k = 4\) row of the earlier \(Y_{2}\)-defining-ideal example (generators omitted here to avoid repeating that display; same growth pattern, quartic with up to 8 terms per generator).
  \end{itemize}
\end{itemize}

\end{example}

\section{Cohen-Macaulayness}\label{sec-cm}

\begin{proposition}{Eagon-Northcott / Hilbert-Burch, applied}{prop-detideal} Let \(M\) be a \(2 \times 3\) matrix with entries in a polynomial ring \(R\) over a field (arbitrary entries, not necessarily independent indeterminates), and \(J\) the ideal of its \(2 \times 2\) minors. If \(J\) has the \textbf{expected} codimension \(3 - 2 + 1 = 2\), then \(J\) is a perfect ideal of grade 2 -- \(\text{pd}(R/J) = 2\), hence \(R/J\) is Cohen-Macaulay.

\end{proposition}

\begin{proof}[{Proof of \thmref{Proposition}{prop-detideal}}] Bruns and Vetter \cite{BrunsVetter1988}, Theorem (2.7), p. 13, gives exactly this for an \textbf{arbitrary} matrix, as needed since one row of our \(2 \times 3\) matrix (\(S\left( {\overline{w}}_{k - 1} \right)\)) is a fixed polynomial vector, not a row of independent indeterminates. (The construction originates with Eagon and Northcott \cite{EagonNorthcott1962}. It is the more general form given by Bruns and Vetter that is applied here.)

\end{proof}

\begin{remark}{why the generator count above wasn't proven this way}{auto5} One might expect \thmref{Proposition}{prop-detideal} to hand over the generator count as well, for free, via Bruns-Vetter's Theorem (16.36), p. 217, the Theorem of Hilbert-Burch. It does not. That theorem as actually stated gives \(I = a \cdot I_{m}(f)\) for a nonzerodivisor \(a\), and a generator count is not among its conclusions.

Nor could it be. At expected codimension the minors of an \textbf{arbitrary} \(2 \times 3\) matrix can be perfect and Cohen-Macaulay while one minor vanishes identically, or is a scalar combination of the other two. A single example settles it: over \(\mathbb{F}\lbrack x,y\rbrack\) the matrix \(\begin{pmatrix}
x & y & 0 \\
0 & 0 & 1
\end{pmatrix}\) has minors \(0,x,y\) and height 2, exactly as expected, yet \(J = (x,y)\) is perfect and Cohen-Macaulay with only \textbf{two} generators.

This is why \thmref{Theorem}{thm-mingens} is proven directly from linear independence rather than cited from this section's machinery.

\end{remark}

\begin{theorem}{\(R_{k}/I_{k}\) is Cohen-Macaulay}{thm-cm} For every \(k \geq 2\): \(R_{k}/I_{k}\) is Cohen-Macaulay.

\end{theorem}

\begin{proof}[{Proof of \thmref{Theorem}{thm-cm}}] By \thmref{Theorem}{thm-codim}, \(I_{k}\) has codimension 2, the expected value for the \(2 \times 3\) matrix of \thmref{Theorem}{thm-determinantal} presenting it. \thmref{Proposition}{prop-detideal} applies directly. (This uses \thmref{Theorem}{thm-codim} only -- not \thmref{Theorem}{thm-mingens}'s generator count, not \thmref{Theorem}{thm-componentcount}'s component count -- genuinely independent of both.)

\end{proof}

\section{Radicality}\label{sec-rad}

Cohen-Macaulayness and the explicit component structure of \cref{sec-geom} combine, via Serre's criterion, to give radicality -- which makes \(I_{k}\) itself, not just \(\sqrt{I_{k}}\), the object carrying the properties established so far, and upgrades \cref{sec-geom}'s component count to a genuine primary decomposition.

Serre's two conditions are written \(\left( R_{n} \right)\) and \(\left( S_{n} \right)\) throughout, parenthesised to keep them apart from this article's polynomial rings \(R_{k}\). A Noetherian ring \(A\) satisfies \(\left( S_{n} \right)\) if \(\text{depth}\left( A_{P} \right) \geq \min(n,\dim A_{P})\) for every prime \(P \subset A\), and \(\left( R_{n} \right)\) if \(A_{P}\) is a \textbf{regular} local ring for every prime \(P\) of height \(\leq n\); in particular \(\left( R_{0} \right)\) asks for regularity at the minimal primes only. Serre's criterion for reducedness \cite{Stacks031R} says that a Noetherian ring is reduced if and only if it satisfies both \(\left( R_{0} \right)\) and \(\left( S_{1} \right)\).

\begin{lemma}{reduction via Serre}{lem-serre} \(R_{k}/I_{k}\) satisfies \(\left( S_{n} \right)\) for every \(n\) -- in particular \(\left( S_{1} \right)\). Consequently \(I_{k}\) is radical if and only if \(R_{k}/I_{k}\) satisfies \(\left( R_{0} \right)\): regularity at each of its minimal primes.

\end{lemma}

\begin{proof}[{Proof of \thmref{Lemma}{lem-serre}}] Cohen-Macaulayness (\thmref{Theorem}{thm-cm}) is the extremal case of \(\left( S_{n} \right)\): it asks for \(\text{depth}\left( A_{P} \right) = \dim A_{P}\) at \textbf{every} prime, so \(\text{depth}\left( A_{P} \right) = \dim A_{P} \geq \min(n,\dim A_{P})\) holds trivially whatever \(n\) is, and a Cohen-Macaulay ring satisfies \(\left( S_{n} \right)\) for every \(n\). (This needs nothing beyond the definitions, and is standard: Bruns and Vetter \cite{BrunsVetter1988}, or Eisenbud \cite{Eisenbud1995}.) Given \(\left( S_{1} \right)\), Serre's criterion makes reducedness of \(R_{k}/I_{k}\) -- radicality of \(I_{k}\) -- equivalent to \(\left( R_{0} \right)\) alone.

\end{proof}

The four lemmas that follow are the standard Jacobian criterion, split into its two genuinely separate ingredients -- one linear-algebraic (\thmref{Lemma}{lem-cotangent}), one dimension-theoretic (\thmref{Lemma}{lem-localdim}) -- before either is used. The order is deliberate: the rank bound and the regularity criterion each draw on \textbf{both} ingredients and on neither each other.

\begin{lemma}{the cotangent identity}{lem-cotangent} Let \(S = \mathbb{F}\left\lbrack x_{1},\ldots,x_{n} \right\rbrack\), \(I = \left( f_{1},\ldots,f_{m} \right)\), \(A = S/I\), and \(x \in V(I)\) an \(\mathbb{F}\)-rational point, with \(\mathfrak{n} \subset S\) the maximal ideal at \(x\) and \(\mathfrak{m} = \mathfrak{n}/I \subset A\) its image. Then \begin{equation}\fitdisplay{\dim_{\mathbb{F}}\mathfrak{m}/\mathfrak{m}^{2} = n - \text{ rank }\text{ Jac}(x).}\end{equation}

\end{lemma}

\begin{proof}[{Proof of \thmref{Lemma}{lem-cotangent}}] Only linear algebra is needed here; in particular, no hypothesis on \(A\) is required. Sending \(g\) to its gradient at \(x\) identifies the cotangent space \(\mathfrak{n}/\mathfrak{n}^{2}\) with \(\mathbb{F}^{n}\). Under that identification the image of \(I\) is the span of the linear parts at \(x\) of \(f_{1},\ldots,f_{m}\): a general element of \(I\) is \(\sum_{i}g_{i}f_{i}\), whose linear part at \(x\) is \(\sum_{i}g_{i}(x)\) times that of \(f_{i}\), since every \(f_{i}\) vanishes at \(x\). That span is the row space of \(\text{Jac}(x)\), of dimension \(\text{rank }\text{ Jac}(x)\). As \(\mathfrak{m}/\mathfrak{m}^{2} \cong \left( \mathfrak{n}/\mathfrak{n}^{2} \right)/\text{ im}(I)\), the claim follows.

\end{proof}

\begin{lemma}{the local dimension at a rational point}{lem-localdim} Let \(S = \mathbb{F}\left\lbrack x_{1},\ldots,x_{n} \right\rbrack\), \(I = \left( f_{1},\ldots,f_{m} \right)\), \(A = S/I\) equidimensional of dimension \(n - c\), and \(x \in V(I)\) an \(\mathbb{F}\)-rational point with maximal ideal \(\mathfrak{m} \subset A\). Write \(\text{Jac}(x) \in \mathbb{F}^{m \times n}\) for the Jacobian \(\left( \partial f_{i}/\partial x_{j} \right)\) \textbf{evaluated at \(x\)}: a matrix of scalars, not the matrix of polynomials it is evaluated from. (The polynomial Jacobian itself appears later as \(J_{k}\) in \thmref{Lemma}{lem-jacrecursion}; the two are not the same object, and only the evaluated one is used in this section.) Then \(\dim A_{\mathfrak{m}} = n - c\).

\end{lemma}

\begin{proof}[{Proof of \thmref{Lemma}{lem-localdim}}] The local dimension is computed componentwise: \begin{equation}\fitdisplay{\dim A_{\mathfrak{m}} = \max\left\{ \dim(A/P)_{\mathfrak{m}}:P \in \text{ Min}(A),P \subseteq \mathfrak{m} \right\},}\end{equation} and each \(A/P = S/\widetilde{P}\) is an affine domain over a field, for which \(\dim\left( S/\widetilde{P} \right)_{\mathfrak{M}} = \dim S/\widetilde{P}\) at \textbf{every} maximal ideal \(\mathfrak{M}\) -- affine domains over a field are catenary and equicodimensional (Eisenbud \cite{Eisenbud1995}, \emph{Commutative Algebra}, Cor. 13.4; equivalently the Stacks Project, tag 00OS \cite{Stacks00OS} -- a dimension formula for finitely generated field extensions, holding over an arbitrary field, not needing a perfect-field hypothesis by accident). By equidimensionality every such \(\dim S/\widetilde{P}\) equals \(n - c\).

\end{proof}

\begin{lemma}{Jacobian criterion: the rank bound}{lem-jacobian-rankbound} The hypotheses are those of \thmref{Lemma}{lem-localdim}. Then \(\text{rank }\text{ Jac}(x) \leq c\).

\end{lemma}

\begin{proof}[{Proof of \thmref{Lemma}{lem-jacobian-rankbound}}] In any Noetherian local ring one has \(\dim_{\mathbb{F}}\mathfrak{m}/\mathfrak{m}^{2} \geq \dim A_{\mathfrak{m}}\). By \thmref{Lemma}{lem-cotangent} the left-hand side is \(n - \text{ rank }\text{ Jac}(x)\), and by \thmref{Lemma}{lem-localdim} the right-hand side is \(n - c\); rearranging gives \(\text{rank }\text{ Jac}(x) \leq c\).

\end{proof}

\begin{lemma}{Jacobian criterion: regularity at rank \(= c\)}{lem-jacobian-regularity} The hypotheses are those of \thmref{Lemma}{lem-localdim}. If \(\text{rank }\text{ Jac}(x) = c\), then \(A_{\mathfrak{m}}\) is regular of dimension \(n - c\); in particular \(x\) lies on a \textbf{unique} component of \(V(I)\), whose local ring there is a field.

\end{lemma}

\begin{proof}[{Proof of \thmref{Lemma}{lem-jacobian-regularity}}] If \(\text{rank }\text{ Jac}(x) = c\) then \(\dim_{\mathbb{F}}\mathfrak{m}/\mathfrak{m}^{2} = n - c\) by \thmref{Lemma}{lem-cotangent}, while \(\dim A_{\mathfrak{m}} = n - c\) by \thmref{Lemma}{lem-localdim}; equality of the two is the definition of regularity. A regular local ring is a domain, so \(\mathfrak{m}\) contains a unique minimal prime of \(A\), and the local ring at that prime is a further localization of a regular local ring at its own maximal ideal in dimension 0, hence a field.

\end{proof}

\begin{lemma}{the Jacobian's own recursion}{lem-jacrecursion} Write \(J_{k}\) for the \(3 \times 3k\) Jacobian of \({\overline{w}}_{k}\), \(\lbrack u\rbrack_{\times}\) for the skew matrix with \(\lbrack u\rbrack_{\times}v = u \times v\). Then \(J_{2} = \left\lbrack - \left\lbrack {\overline{v}}_{2} \right\rbrack_{\times}|\left\lbrack {\overline{v}}_{1} \right\rbrack_{\times} \right\rbrack\), and for \(k \geq 3\), \begin{equation}\fitdisplay{J_{k} = \left\lbrack - \left\lbrack {\overline{v}}_{k} \right\rbrack_{\times}SJ_{k - 1}|\left\lbrack {S\left( {\overline{w}}_{k - 1} \right)} \right\rbrack_{\times} \right\rbrack.}\end{equation}

\end{lemma}

\begin{proof}[{Proof of \thmref{Lemma}{lem-jacrecursion}}] For the second block, \({\overline{w}}_{k} = {S\left( {\overline{w}}_{k - 1} \right)} \times {\overline{v}}_{k} = \left\lbrack {S\left( {\overline{w}}_{k - 1} \right)} \right\rbrack_{\times}{\overline{v}}_{k}\) is linear in \({\overline{v}}_{k}\) with coefficient matrix free of \({\overline{v}}_{k}\). For the first, differentiating \({\overline{w}}_{k} = {S\left( {\overline{w}}_{k - 1} \right)} \times {\overline{v}}_{k}\) with respect to any variable \(t\) among \({\overline{v}}_{1},\ldots,{\overline{v}}_{k - 1}\) gives \(\partial_{t}{\overline{w}}_{k} = \left( S\partial_{t}{\overline{w}}_{k - 1} \right) \times {\overline{v}}_{k} = - \left\lbrack {\overline{v}}_{k} \right\rbrack_{\times}S\partial_{t}{\overline{w}}_{k - 1}\), using \(a \times b = - \lbrack b\rbrack_{\times}a\); assembling columns gives the stated product.

\end{proof}

\begin{example}{\(J_{3}\), explicitly}{auto6} With variables ordered \(v_{1,1}\), \(v_{1,2}\), \(v_{1,3}\), \(v_{2,1}\), \(v_{2,2}\), \(v_{2,3}\), \(v_{3,1}\), \(v_{3,2}\), \(v_{3,3}\), the \(3 \times 9\) Jacobian of \({\overline{w}}_{3}\) (the three generators of \(I_{3}\) from the earlier example), column-blocked to match \thmref{Lemma}{lem-jacrecursion}'s formula \(J_{3} = \left\lbrack - \left\lbrack {\overline{v}}_{3} \right\rbrack_{\times}SJ_{2}|\left\lbrack {S\left( {\overline{w}}_{2} \right)} \right\rbrack_{\times} \right\rbrack\) -- shown here as the two blocks separately, since side by side the full \(3 \times 9\) matrix does not fit the page width. First, \(- \left\lbrack {\overline{v}}_{3} \right\rbrack_{\times}SJ_{2}\) (\(3 \times 6\)): \begin{equation}\fitdisplay{\begin{pmatrix}
 - v_{2,2}v_{3,2} + v_{2,3}v_{3,3} & v_{2,1}v_{3,2} & - v_{2,1}v_{3,3} & v_{1,2}v_{3,2} - v_{1,3}v_{3,3} & - v_{1,1}v_{3,2} & v_{1,1}v_{3,3} \\
v_{2,2}v_{3,1} & - v_{2,1}v_{3,1} + v_{2,3}v_{3,3} & - v_{2,2}v_{3,3} & - v_{1,2}v_{3,1} & v_{1,1}v_{3,1} - v_{1,3}v_{3,3} & v_{1,2}v_{3,3} \\
 - v_{2,3}v_{3,1} & - v_{2,3}v_{3,2} & v_{2,1}v_{3,1} + v_{2,2}v_{3,2} & v_{1,3}v_{3,1} & v_{1,3}v_{3,2} & - v_{1,1}v_{3,1} - v_{1,2}v_{3,2}
\end{pmatrix}}\end{equation} and then \(\left\lbrack {S\left( {\overline{w}}_{2} \right)} \right\rbrack_{\times}\) (\(3 \times 3\), the skew matrix of \(S\left( {\overline{w}}_{2} \right)\)'s three coordinates from the earlier matrix example): \begin{equation}\fitdisplay{\begin{pmatrix}
0 & v_{1,2}v_{2,1} - v_{1,1}v_{2,2} & - v_{1,3}v_{2,1} + v_{1,1}v_{2,3} \\
 - v_{1,2}v_{2,1} + v_{1,1}v_{2,2} & 0 & - v_{1,3}v_{2,2} + v_{1,2}v_{2,3} \\
v_{1,3}v_{2,1} - v_{1,1}v_{2,3} & v_{1,3}v_{2,2} - v_{1,2}v_{2,3} & 0
\end{pmatrix}}\end{equation} Both blocks computed directly, not merely asserted from the recursive formula.

\end{example}

\begin{remark}{a false shortcut}{auto7} A naive reading of \thmref{Lemma}{lem-jacrecursion} suggests appending \({\overline{v}}_{k} = 0\) preserves rank on an inherited component -- this is false: the old-variable block is left-multiplied by \(- \left\lbrack {\overline{v}}_{k} \right\rbrack_{\times}S\), not merely padded with zero rows. At the first inherited case (\(k = 2 \rightarrow 3\)), appending \({\overline{v}}_{3} = 0\) actually drops the rank to 0. The correct witness (\thmref{Lemma}{lem-witness-inherited} below) instead picks \({\overline{v}}_{k} = e_{i}\) for a standard basis vector \textbf{not} lying in \(S \cdot \text{ im }J_{k - 1}(q)\) -- a genuine choice, not a formality.

\end{remark}

\begin{lemma}{rank-2 witness on the new component}{lem-witness-new} For every \(k \geq 3\), there is a point \(x = (p,0)\), \(p\) the \thmref{Lemma}{lem-nondeg} non-degeneracy witness at level \(k - 1\), with \(x \in V\left( I_{k} \right)\) and \(\text{rank }J_{k(x)} = 2\). At \(k = 2\) the same conclusion holds with \(x = (p,0)\) for any \(p \neq 0\).

\end{lemma}

\begin{proof}[{Proof of \thmref{Lemma}{lem-witness-new}}] \textbf{The case \(k = 2\) separately}, because the statement above cannot be read at \(k = 2\): it would refer to a non-degeneracy witness ``at level 1'', and neither \({\overline{w}}_{1}\) nor \thmref{Lemma}{lem-nondeg} exists there, since \thmref{Definition}{def-wk} starts the recursion at \(j = 2\). Nor does \thmref{Lemma}{lem-jacrecursion}'s \(k \geq 3\) formula apply. Argue directly instead: \(J_{2} = \left\lbrack - \left\lbrack {\overline{v}}_{2} \right\rbrack_{\times}|\left\lbrack {\overline{v}}_{1} \right\rbrack_{\times} \right\rbrack\) carries no \(S\) at all, so for any \(p \neq 0\), \(J_{2}(p,0) = \left\lbrack 0|\lbrack p\rbrack_{\times} \right\rbrack\), which has rank 2 by the kernel computation below. And \((p,0) \in V\left( I_{2} \right)\) since \({\overline{w}}_{2}(p,0) = p \times 0 = 0\).

\textbf{For \(k \geq 3.\)} By \thmref{Lemma}{lem-jacrecursion}, \(J_{k(p,0)} = \left\lbrack 0|\left\lbrack {S\left( {\overline{w}}_{k - 1}(p) \right)} \right\rbrack_{\times} \right\rbrack\) since \(- \lbrack 0\rbrack_{\times}SJ_{k - 1}(p) = 0\); and \({S\left( {\overline{w}}_{k - 1}(p) \right)} \neq 0\) since \({\overline{w}}_{k - 1}(p) \neq 0\) (\thmref{Lemma}{lem-nondeg}) and \(S\) is invertible. For \(u \neq 0\), \(\ker\lbrack u\rbrack_{\times} = \left\{ v:u \times v = 0 \right\} = \mathbb{F}u\) (a cross product vanishes exactly on linearly dependent pairs), so \(\left\lbrack {S\left( {\overline{w}}_{k - 1}(p) \right)} \right\rbrack_{\times}\) has rank 2, giving \(\text{rank }J_{k(p,0)} = 2\).

\end{proof}

\begin{lemma}{rank-2 witness on an inherited component}{lem-witness-inherited} For every \(k \geq 3\) and every component \(Z\) of \(V\left( I_{k - 1} \right)\) with a rank-2 witness \(q \in Z\) (\thmref{Lemma}{lem-witness-new}, or this lemma applied at \(k - 1\)), there is \(x = \left( q,e_{i} \right) \in Z \times \mathbb{A}^{3} \subset V\left( I_{k} \right)\) with \(\text{rank }J_{k(x)} = 2\).

\end{lemma}

\begin{proof}[{Proof of \thmref{Lemma}{lem-witness-inherited}}] Let \(W \coloneqq  S \cdot \text{ im }J_{k - 1}(q)\), a 2-dimensional subspace of \(\mathbb{F}^{3}\) since \(S\) is invertible and \(\text{rank }J_{k - 1}(q) = 2\). Since \(e_{1},e_{2},e_{3}\) span \(\mathbb{F}^{3}\) and \(W \neq \mathbb{F}^{3}\), some \(e_{i} \notin W\); fix such an \(i\). Since \(q \in V\left( I_{k - 1} \right)\), \({\overline{w}}_{k - 1}(q) = 0\), so \thmref{Lemma}{lem-jacrecursion} gives \(J_{k\left( q,e_{i} \right)} = \left\lbrack - \left\lbrack e_{i} \right\rbrack_{\times}SJ_{k - 1}(q)|0 \right\rbrack\), of rank \(\dim(\left\lbrack e_{i} \right\rbrack_{\times}W)\) since \(SJ_{k - 1}(q)\) has image exactly \(W\). As \(\ker\left\lbrack e_{i} \right\rbrack_{\times} = \mathbb{F}e_{i}\) and \(W \cap \mathbb{F}e_{i} = 0\) (else \(e_{i} \in W\)), \(\left\lbrack e_{i} \right\rbrack_{\times}\) is injective on \(W\), so \(\dim(\left\lbrack e_{i} \right\rbrack_{\times}W) = 2\).

\end{proof}

\begin{theorem}{\(I_{k}\) is radical}{thm-radical} For every \(k \geq 2\): \(I_{k}\) is radical.

\end{theorem}

\begin{proof}[{Proof of \thmref{Theorem}{thm-radical}}] By induction on \(k\) (mirroring \thmref{Theorem}{thm-componentcount}'s own induction), \thmref{Lemma}{lem-witness-new} and \thmref{Lemma}{lem-witness-inherited} exhibit a rank-2-Jacobian point on \textbf{every} one of the \(k - 1\) components of \thmref{Theorem}{thm-componentcount}. \thmref{Lemma}{lem-jacobian-regularity} then gives \(\left( R_{0} \right)\) at every one (and, by its ``unique component'' clause, confirms each witness lies only on its intended component -- an independent check on \thmref{Theorem}{thm-componentcount}'s count, not merely assumed); \thmref{Lemma}{lem-serre} gives \(\left( S_{1} \right)\), using \thmref{Theorem}{thm-cm}'s Cohen-Macaulayness. Serre's criterion \cite{Stacks031R} gives radicality.

\end{proof}

This has been verified independently in Sage, for the signed encoding, at every \(k = 2,\ldots,6\): the Jacobian recursion (\thmref{Lemma}{lem-jacrecursion}) checked as a polynomial identity, not just at points, and every witness confirmed to lie on \(V\left( I_{k} \right)\) with \(\text{rank }J_{k} = 2\) (code/referee\_radicality\_S\_check.sage). The sharper claim, that each witness lies on \textbf{exactly} its intended component, is confirmed by an independently-computed primary decomposition at every \(k = 3,4,5,6\), which returns \(k - 1\) components in each case, in agreement with \thmref{Theorem}{thm-componentcount}; the \(k = 6\) decomposition takes some hours.

\section{The Primary Decomposition Theorem}

Every ingredient is now in hand for the fully general statement -- component count (\cref{sec-geom}) and radicality (\cref{sec-rad}) together, neither sufficient alone. One convention has to be made explicit first, because it is the hinge on which the statement's claimed generality turns.

\begin{remark}{the geometric arguments run over \(\overline{\mathbb{F}}\), and the conclusion descends}{rem-base-field} The three preceding sections argue geometrically: they speak of \(V\left( I_{k} \right)\), of its irreducible components, and of the closure of the image of a parametrization, and at the decisive step they identify \(\text{Min}\left( I_{k} \right)\) with the vanishing ideals of those components. That identification is a Nullstellensatz statement and is \textbf{false over a general field} -- over \(\mathbb{F}_{3}\), for instance, \(V\left( I_{k} \right)\left( \mathbb{F}_{3} \right)\) is a finite set of points and carries no such information at all. So throughout those sections, and in the theorem below, \(V( \cdot )\) means the vanishing locus in \(\mathbb{A}^{3k}\left( \overline{\mathbb{F}} \right)\) over a fixed algebraic closure, and ``component'' means an irreducible component there. This is the usual convention and is what the classical literature on the commuting variety assumes as well (Artin-Hochster is stated over an algebraically closed field in \cite{MajidiZolbaninSnapp}).

The conclusion nonetheless descends to \(\mathbb{F}\) itself, so the theorem may be read over any field of characteristic \(\neq 2\), as stated. The reason is that each \(P_{j}\) is not merely prime but \textbf{absolutely} prime. By \thmref{Definition}{def-primecomponents}, \(P_{j}\) \textbf{is} the kernel of an \(\mathbb{F}\)-algebra homomorphism \(\varphi_{j}:R_{k} \rightarrow \mathbb{F}\left\lbrack t_{1},\ldots,t_{3k - 2} \right\rbrack\) into a polynomial ring, so \(R_{k}/P_{j}\) embeds in that polynomial ring. Since \(\overline{\mathbb{F}}\) is flat over \(\mathbb{F}\), tensoring preserves the embedding: \begin{equation}\fitdisplay{\overline{\mathbb{F}} \otimes_{\mathbb{F}}\left( R_{k}/P_{j} \right) \hookrightarrow \overline{\mathbb{F}}\left\lbrack t_{1},\ldots,t_{m} \right\rbrack,}\end{equation} whose target is a domain, so the extension of \(P_{j}\) to \(\overline{\mathbb{F}} \otimes_{\mathbb{F}}R_{k}\) is again prime. Formation of a finite intersection of ideals likewise commutes with the flat base change \(\mathbb{F} \rightarrow \overline{\mathbb{F}}\). Hence the displayed decomposition, established over \(\overline{\mathbb{F}}\), holds already over \(\mathbb{F}\), with the same \(k - 1\) primes.

\end{remark}

\begin{theorem}{Primary Decomposition Theorem}{thm-primarydecomp} For every \(k \geq 2\), over any field of characteristic \(\neq 2\): \(I_{k}\) has exactly \(k - 1\) irreducible, radical, primary components, namely \(P_{2},\ldots,P_{k}\) of \thmref{Definition}{def-primecomponents}, and \begin{equation}\fitdisplay{I_{k} = P_{2} \cap \cdots \cap P_{k}}\end{equation} is an irredundant primary decomposition into \(k - 1\) primes.

\end{theorem}

\begin{proof}[{Proof of \thmref{Theorem}{thm-primarydecomp}}] The two halves were proven independently above, so this combines them in full rather than asserting the conjunction in one line.

\textbf{The decomposition itself.} By \thmref{Theorem}{thm-radical}, \(I_{k}\) is radical, so \(I_{k} = \sqrt{I_{k}}\). A radical ideal in a Noetherian ring equals the intersection of its own minimal primes: \begin{equation}\fitdisplay{I_{k} = \cap_{P \in \text{ Min}\left( I_{k} \right)}P.}\end{equation} This holds because \(\sqrt{I_{k}}\) is the intersection of \textbf{all} primes containing \(I_{k}\) -- Krull's theorem, not the definition of the radical, which is instead the set of elements with a power in \(I_{k}\) -- and the minimal primes already suffice. Every non-minimal prime containing \(I_{k}\) contains a minimal one, and is therefore redundant in the intersection.

\thmref{Theorem}{thm-componentcount} identifies \(\text{Min}\left( I_{k} \right)\) explicitly, over \(\overline{\mathbb{F}}\) as \thmref{Remark}{rem-base-field} requires: exactly \(k - 1\) primes, one per block pattern \(j = 2,\ldots,k\), named \(P_{2},\ldots,P_{k}\) in \thmref{Definition}{def-primecomponents}. Substituting them gives the displayed equation.

\textbf{Why the components are primary.} Each \(P_{j}\) is prime, so it is trivially its own primary-decomposition component, being \(P_{j}\)-primary with associated prime \(P_{j}\). Nothing beyond primality has to be argued.

\textbf{Why they are prime.} Nothing is left to prove here either. \thmref{Definition}{def-primecomponents} \textbf{defines} \(P_{j}\) as \(\ker\varphi_{j}\), and the kernel of a homomorphism into a domain is prime. That argument covers \(j = 2\) along with the rest, at no extra cost. \thmref{Lemma}{lem-y2-irred} would give the same conclusion for \(j \geq 3\), but it is stated for \(k \geq 3\) and so misses \(j = 2\); \(P_{2}\) is \(I_{2}\) itself, the ideal of maximal minors of a generic \(2 \times 3\) matrix, which is classically prime. Neither of those routes is needed.

Extending \(P_{j}\) from \(R_{j}\) to \(R_{k}\) by adjoining the free blocks \({\overline{v}}_{j + 1},\ldots,{\overline{v}}_{k}\) preserves primality, since \(R_{k}/P_{j}R_{k}\) is a polynomial ring over the domain \(R_{j}/P_{j}\).

\textbf{Irredundancy.} No \(P_{j}\) can be omitted, i.e. no \(P_{j}\) contains \(\cap_{i \neq j}P_{i}\). This is exactly \thmref{Theorem}{thm-componentcount}'s own no-containment conclusion, obtained from \thmref{Proposition}{prop-y2-not-in-y1} and \thmref{Proposition}{prop-y1-not-in-y2} applied at every level of that theorem's induction. It needs no further proof here.

\end{proof}

\begin{remark}{summary}{auto8} \thmref{Theorem}{thm-codim} (codimension 2), \thmref{Theorem}{thm-mingens} (3 minimal generators), \thmref{Theorem}{thm-cm} (Cohen-Macaulay), and \thmref{Theorem}{thm-radical} (radical) together say exactly what a single bundled ``\(I_{k}\) has codimension 2, is radical, has exactly 3 minimal generators, and is Cohen-Macaulay'' statement would -- now as four independently-proven results, each established at the point in this article where its actual dependencies (and no others) were available, rather than one claim proven in parts across several sections.

\end{remark}

\section{Summary}

For every \(k \geq 2\), over any field \(\mathbb{F}\) of characteristic \(\neq 2\):

\begin{itemize}
\item
  \(I_{k}\) has codimension 2 (\thmref{Theorem}{thm-codim}), presented as the ideal of \(2 \times 2\) minors of an explicit \(2 \times 3\) matrix (\thmref{Theorem}{thm-determinantal}) -- a genuine determinantal ideal, not merely determinantal-looking.
\item
  \(I_{k}\) has exactly 3 minimal generators (\thmref{Theorem}{thm-mingens}), proven directly from linear independence, independent of the determinantal presentation.
\item
  \(R_{k}/I_{k}\) is Cohen-Macaulay (\thmref{Theorem}{thm-cm}), from the determinantal presentation and classical Eagon-Northcott/Hilbert-Burch theory (\thmref{Proposition}{prop-detideal}).
\item
  \(I_{k}\) is radical (\thmref{Theorem}{thm-radical}), via Serre's criterion: Cohen-Macaulayness gives \(\left( S_{1} \right)\) (\thmref{Lemma}{lem-serre}), and an explicit rank-2 Jacobian witness on every component gives \(\left( R_{0} \right)\) (\thmref{Lemma}{lem-witness-new}, \thmref{Lemma}{lem-witness-inherited}).
\item
  \(V\left( I_{k} \right)\) has exactly \(k - 1\) irreducible components (\thmref{Theorem}{thm-componentcount}), each rational -- in fact birational to an explicit affine parameter space (\thmref{Proposition}{prop-rational}). Radicality and this component count together -- rationality plays no part in the assembly -- give the irredundant primary decomposition of the Primary Decomposition Theorem (\thmref{Theorem}{thm-primarydecomp}): \(I_{k} = P_{2} \cap \cdots \cap P_{k}\).
\end{itemize}

\section{Methods}\label{sec-methods}

This article was produced by a workflow that is not the usual one. Since the workflow bears on how some of the claims above should be read, it is set out here rather than left to be inferred.

\textbf{Computer algebra, in two roles that should not be confused.} SageMath and Macaulay2 were used throughout, for two different purposes. The first is \textbf{conjecture}: the component count of \thmref{Theorem}{thm-componentcount}, its block-involvement pattern, and the shape of the primary decomposition were found by computing primary decompositions at small \(k\) and reading off the pattern, well before any of it was proved. The second is \textbf{verification}: once proved, the general statements were checked against independent computation -- the determinantal structure by Gröbner-basis equality in both directions and, separately, against ordinary \(2 \times 2\) matrix multiplication (the two routes sharing no code); the Jacobian recursion as a polynomial identity rather than at sampled points; the witness points for position and rank, on exactly the intended components; and the decomposition itself recomputed at \(k = 3,4,5,6\).

Because rank, dimension and radicality are \textbf{not} characteristic-independent (a matrix of rank 3 over \(\mathbb{Q}\) may drop rank modulo \(p\)), computations over \(\mathbb{Q}\) license nothing about \(\mathbb{F}_{p}\), and the four main results were therefore re-run over \(\mathbb{F}_{2}\), \(\mathbb{F}_{3}\), \(\mathbb{F}_{4}\), \(\mathbb{F}_{5}\), \(\mathbb{F}_{7}\) and \(\mathbb{F}_{9}\) as well. No deviation appeared. The scripts are named at the point of use, and are deposited, with their full output transcripts, at Zenodo (DOI: \href{https://doi.org/10.5281/zenodo.23433730}{10.5281/zenodo.23433730}, with \href{https://doi.org/10.5281/zenodo.23433729}{10.5281/zenodo.23433729} for all versions); the paths cited in the text are the paths inside that archive. The deposit's \texttt{runs/REFERENCE-CERTIFICATE.md} records, for each cited source, the exact statement relied on and how it was verified. The explicit computations are displayed in \cref{sec-appendix}.

Nothing in this article is conjectural, or checked only at small \(k\): every statement is proven for \textbf{every} \(k \geq 2\), over an \textbf{arbitrary} field of the stated characteristic. The computational thread corroborates the general proofs; nowhere does it substitute for one.

\textbf{A partial formalization.} The statements were formalized in Lean 4 against \emph{mathlib} and machine-checked, with no proof gaps and no axioms beyond the three that every \emph{mathlib} development uses. The formalization reaches part of the article, and the division is this.

Formalized:

\begin{itemize}
\item
  the basic identity (\thmref{Theorem}{thm-basic-identity}), in the strengthened form that drops the traceless hypothesis;
\item
  the minimal-generator theorem (\thmref{Theorem}{thm-mingens}) for the actual \({\overline{w}}_{k}\), with the chain beneath it: that independent same-degree generators are minimal (\thmref{Lemma}{lem-mingens-generic}), the homogeneity of \({\overline{w}}_{k}\), the freshness of each new variable block, and the linear independence itself (\thmref{Lemma}{lem-mingens-indep});
\item
  the Jacobian recursion (\thmref{Lemma}{lem-jacrecursion});
\item
  assorted smaller computational claims of \cref{sec-setup} and \cref{sec-geom}.
\end{itemize}

Not formalized, the libraries for it not presently existing:

\begin{itemize}
\item
  the codimension (\thmref{Theorem}{thm-codim});
\item
  Cohen-Macaulayness (\thmref{Theorem}{thm-cm});
\item
  radicality (\thmref{Theorem}{thm-radical});
\item
  the component count (\thmref{Theorem}{thm-componentcount});
\item
  the Primary Decomposition Theorem (\thmref{Theorem}{thm-primarydecomp}).
\end{itemize}

Everything in the second list rests on determinantal-ideal theory or on dimension theory. The guarantee is therefore partial, and does not cover the article's main results.

\section{Acknowledgements}

The author first learned of the ring of commuting matrices some thirty years ago from Professor Jan-Erik Roos, in a course he held on the then fledgling field of computer algebra. It is a pleasant debt to record here, given how much of what follows was found with a computer algebra system open beside the page.

The computations were carried out in SageMath and Macaulay2, and the formalization in Lean 4 with \emph{mathlib}; this article would have a different and smaller content without them. The nature and extent of the assistance from a large language model is set out in \cref{sec-methods} and in the declaration at the end.

\section{Appendix: the explicit computations}\label{sec-appendix}

The results above are stated with their computations named rather than displayed, so that the argument stays legible. This appendix displays them. Everything here is \textbf{corroboration}: no proof in the article rests on a computation, and nothing below is offered as one. Each block names the script that produced it; the scripts and their full transcripts are in the Zenodo deposit (DOI: \href{https://doi.org/10.5281/zenodo.23433730}{10.5281/zenodo.23433730}).

(The scripts write \texttt{D} internally for the matrix this article calls \(S\) -- the same matrix, \(\text{diag}( - 1, - 1,1)\).)

\subsection{Both sides of the basic identity}\label{app-identity}

\thmref{Theorem}{thm-basic-identity} asserts \(\text{enc}\left( \left\lbrack A_{0},B_{0} \right\rbrack \right) = {S\left( \text{enc}\left( A_{0} \right) \times \text{ enc}\left( B_{0} \right) \right)}\). Expanded over generic entries \(a_{11},\ldots,b_{22}\), the two sides are, component by component:

\begin{verbatim}
lhs_1 = rhs_1 = -2*a21*b12 + 2*a12*b21
lhs_2 = rhs_2 = -a12*b11 + a21*b11 + a11*b12 - a22*b12 - a11*b21 + a22*b21
                + a12*b22 - a21*b22
lhs_3 = rhs_3 = -a12*b11 - a21*b11 + a11*b12 - a22*b12 + a11*b21 - a22*b21
                + a12*b22 + a21*b22
\end{verbatim}

The difference is \((0,0,0)\) identically. Chaining once more gives the triple commutator \(\text{enc}\left( \left\lbrack \lbrack A,B\rbrack,C \right\rbrack \right)\), which the recursion computes as \(S\left( {S\left( \text{enc}(A) \times \text{ enc}(B) \right)} \times \text{ enc}(C) \right)\); its three components have \(8\), \(20\) and \(20\) terms:

\begin{verbatim}
comp_1 = -2*a21*b11*c12 + 2*a11*b21*c12 - 2*a22*b21*c12 + 2*a21*b22*c12
         - 2*a12*b11*c21 + 2*a11*b12*c21 - 2*a22*b12*c21 + 2*a12*b22*c21

comp_2 = a12*b11*c11 + a21*b11*c11 - a11*b12*c11 + a22*b12*c11 - a11*b21*c11
         + a22*b21*c11 - a12*b22*c11 - a21*b22*c11 - 2*a21*b12*c12
         + 2*a12*b21*c12 + 2*a21*b12*c21 - 2*a12*b21*c21 - a12*b11*c22
         - a21*b11*c22 + a11*b12*c22 - a22*b12*c22 + a11*b21*c22
         - a22*b21*c22 + a12*b22*c22 + a21*b22*c22

comp_3 = a12*b11*c11 - a21*b11*c11 - a11*b12*c11 + a22*b12*c11 + a11*b21*c11
         - a22*b21*c11 - a12*b22*c11 + a21*b22*c11 - 2*a21*b12*c12
         + 2*a12*b21*c12 - 2*a21*b12*c21 + 2*a12*b21*c21 - a12*b11*c22
         + a21*b11*c22 + a11*b12*c22 - a22*b12*c22 - a11*b21*c22
         + a22*b21*c22 + a12*b22*c22 - a21*b22*c22
\end{verbatim}

The factors of \(2\) here are expected: on the traceless part \(a_{22} = - a_{11}\), so \(d_{a} = a_{11} - a_{22} = 2a_{11}\). The \(\pm 1\)-coefficient claim of \thmref{Example}{ex-wk-signed} concerns the \(\overline{v}\)-coordinates, a different expansion in different variables. Both routes agree exactly (code/regenerate\_displays\_S.sage, code/verify\_lorentzian\_encoding\_identity.sage).

\subsection{The sign twist leaves the ideal unchanged}\label{app-remark2}

The three components of \({\overline{w}}_{2}\), which are also the entries of \(S\left( {\overline{w}}_{2} \right)\) up to sign, are \begin{equation}\fitdisplay{- v_{1,3}v_{2,2} + v_{1,2}v_{2,3},\quad v_{1,3}v_{2,1} - v_{1,1}v_{2,3},\quad - v_{1,2}v_{2,1} + v_{1,1}v_{2,2},}\end{equation} so \(\text{ideal}\left( {\overline{w}}_{2} \right) = \text{ ideal}\left( S{\overline{w}}_{2} \right)\) generator by generator: two of the three are negated and the third is fixed, and negating a generator does not change the ideal it generates -- this is \thmref{Definition}{def-wk}'s ``the two triples generate the same ideal'', displayed. Confirmed independently by Gröbner-basis equality in code/verify\_encoding\_gl3\_equivariance.sage. Incidentally \(S^{2} = I\) exactly, so the twist is an involution.

\subsection{\texorpdfstring{The generators of \(I_{6}\)}{The generators of I\_\{6\}}}\label{app-i6}

\thmref{Theorem}{thm-mingens} gives \(I_{k}\) three minimal generators for every \(k\). At \(k = 6\) they are degree-\(6\) forms in the \(18\) variables \(v_{1,1},\ldots,v_{6,3}\), with \(32\) terms each and \textbf{every coefficient equal to \(\pm 1\)}:

\begin{verbatim}
w_6,1 = -v1_3*v2_1*v3_1*v4_1*v5_1*v6_2 + v1_1*v2_3*v3_1*v4_1*v5_1*v6_2
        - v1_3*v2_2*v3_2*v4_1*v5_1*v6_2 + v1_2*v2_3*v3_2*v4_1*v5_1*v6_2
        - v1_2*v2_1*v3_2*v4_3*v5_1*v6_2 + v1_1*v2_2*v3_2*v4_3*v5_1*v6_2
        + v1_3*v2_1*v3_3*v4_3*v5_1*v6_2 - v1_1*v2_3*v3_3*v4_3*v5_1*v6_2
        - v1_3*v2_1*v3_1*v4_2*v5_2*v6_2 + v1_1*v2_3*v3_1*v4_2*v5_2*v6_2
        - v1_3*v2_2*v3_2*v4_2*v5_2*v6_2 + v1_2*v2_3*v3_2*v4_2*v5_2*v6_2
        + v1_2*v2_1*v3_1*v4_3*v5_2*v6_2 - v1_1*v2_2*v3_1*v4_3*v5_2*v6_2
        + v1_3*v2_2*v3_3*v4_3*v5_2*v6_2 - v1_2*v2_3*v3_3*v4_3*v5_2*v6_2
        + v1_2*v2_1*v3_1*v4_1*v5_1*v6_3 - v1_1*v2_2*v3_1*v4_1*v5_1*v6_3
        + v1_3*v2_2*v3_3*v4_1*v5_1*v6_3 - v1_2*v2_3*v3_3*v4_1*v5_1*v6_3
        + v1_2*v2_1*v3_2*v4_2*v5_1*v6_3 - v1_1*v2_2*v3_2*v4_2*v5_1*v6_3
        - v1_3*v2_1*v3_3*v4_2*v5_1*v6_3 + v1_1*v2_3*v3_3*v4_2*v5_1*v6_3
        + v1_3*v2_1*v3_1*v4_2*v5_3*v6_3 - v1_1*v2_3*v3_1*v4_2*v5_3*v6_3
        + v1_3*v2_2*v3_2*v4_2*v5_3*v6_3 - v1_2*v2_3*v3_2*v4_2*v5_3*v6_3
        - v1_2*v2_1*v3_1*v4_3*v5_3*v6_3 + v1_1*v2_2*v3_1*v4_3*v5_3*v6_3
        - v1_3*v2_2*v3_3*v4_3*v5_3*v6_3 + v1_2*v2_3*v3_3*v4_3*v5_3*v6_3
\end{verbatim}

The other two components, \(w_{6,2}\) and \(w_{6,3}\), have the same shape and the same \(32\)-term, \(\pm 1\) profile; all three are printed in full in code/regenerate\_displays\_S.sage's transcript, which ships as \texttt{runs/display-ground-truth-S.txt}. The script asserts the coefficient profile rather than merely printing it.

\subsection{The non-degeneracy witness}\label{app-witness}

The following table evaluates the point constructed in the proof of \thmref{Lemma}{lem-nondeg}.

\begin{verbatim}
j    witness (v_1,...,v_j)              w_j
2    (e1,e2)                            e3
3    (e1,e2,e1)                         e2
4    (e1,e2,e1,e1)                      e3
5    (e1,e2,e1,e1,e1)                   e2
6    (e1,e2,e1,e1,e1,e1)                e3
7    (e1,e2,e1,e1,e1,e1,e1)             e2
8    (e1,e2,e1,e1,e1,e1,e1,e1)          e3
\end{verbatim}

Each level is checked twice in code/verify\_nondegeneracy\_witness.sage: once as the induction runs, and once by feeding the finished witness back through \thmref{Definition}{def-wk}'s recursion as a point of \(\mathbb{A}^{3j}\), which is how \thmref{Lemma}{lem-nondeg} states it.

\subsection{Independence of the generators}\label{app-independence}

This is the corroboration cited in the proof of \thmref{Lemma}{lem-mingens-indep}. For each \(k\), the three components of \({\overline{w}}_{k}\) are expanded over the monomials that actually occur and the resulting \(3 \times N\) coefficient matrix is checked to have rank \(3\), which is exactly the assertion that no nontrivial constant relation holds among them:

\begin{longtable}[]{@{}cccccc@{}}
\toprule\noalign{}
\endhead
\bottomrule\noalign{}
\endlastfoot
\(k\) & \#variables & degree & \#monomials & rank & independent? \\
2 & 6 & 2 & 6 & 3 & yes \\
3 & 9 & 3 & 12 & 3 & yes \\
4 & 12 & 4 & 24 & 3 & yes \\
5 & 15 & 5 & 48 & 3 & yes \\
6 & 18 & 6 & 96 & 3 & yes \\
\end{longtable}

The monomial count doubles at each step while the rank stays at \(3\), so the margin widens rather than narrows as \(k\) grows (code/verify\_generator\_independence.sage).

\section{Declaration of generative AI and AI-assisted technologies in the manuscript preparation process}

During the preparation of this work the author used Claude (Anthropic), through Claude Code, in order to assist with mathematical exploration, drafting, computational verification and the preparation of the Lean formalization. The models used were Claude Sonnet 5 (August 2026), Claude Opus 5 (August to October 2026) and Claude Opus 5.5 (October 2026), with adversarial reviews of the proofs run with Claude Fable 5 (August and September 2026). After using this tool, the author reviewed and edited the content as needed and takes full responsibility for the content of the published article.

The assistance was substantial, and its extent is stated here rather than left to be inferred. The itemised account is \cref{sec-methods}: which results were machine-checked and which were not, what the computer algebra established and what it merely corroborated, and what the adversarial review found. Claude is not a co-author, and is not credited as one.

Every computational claim has been checked along a route independent of the statement being tested, and the reference list has been checked against the sources cited, with the one exception that where a source is a printed book outside this project's open-access policy, an independently verified equivalent is cited alongside it. The author has verified the results to the best of his ability, made the final edits, and assumes full responsibility for any errors, mistakes or gaps that remain.

\section{References}

\begin{thebibliography}{99}
\bibitem{BrunsVetter1988} W. Bruns, U. Vetter, \emph{Determinantal Rings}, Lecture Notes in Mathematics 1327, Springer-Verlag, 1988.  \bibitem{EagonNorthcott1962} J.A. Eagon, D.G. Northcott, \emph{Ideals defined by matrices and a certain complex associated with them}, Proc. Roy. Soc. London Ser. A 269(1337) (1962), 188-204. DOI: \href{https://doi.org/10.1098/rspa.1962.0170}{10.1098/rspa.1962.0170}.  \bibitem{MotzkinTaussky1955} T. Motzkin, O. Taussky, \emph{Pairs of matrices with property L. II}, Trans. Amer. Math. Soc. 80 (1955), 387-401.  \bibitem{Stacks00OS} \emph{The Stacks Project}, tag 00OS (dimension formula for finitely generated field extensions / equicodimensionality of affine domains), \href{https://stacks.math.columbia.edu/tag/00OS}{stacks.math.columbia.edu/tag/00OS}.  \bibitem{Gerstenhaber1961} M. Gerstenhaber, \emph{On dominance and varieties of commuting matrices}, Ann. of Math. 73 (1961), 324-348.  \bibitem{ArtinHochsterConjecture} M. Artin, M. Hochster, attributed conjecture: for \(I\) the ideal generated by the entries of \(XY - YX\), \(\text{Spec}(R/I)\) is reduced and Cohen-Macaulay. No single dateable original publication found; stated as ``Conjecture 1.1 (M. Artin, M. Hochster)'' in Majidi-Zolbanin and Snapp below.  \bibitem{MajidiZolbaninSnapp} M. Majidi-Zolbanin, B. Snapp, \emph{A note on the variety of pairs of matrices whose product is symmetric}, in \emph{Commutative Algebra and Its Connections to Geometry}, Contemporary Mathematics 555, American Mathematical Society, 2011, pp. 145-150.  \bibitem{Chau2024} T. Chau, \emph{The \(F\)-regularity of algebraic sets related to commutator matrices}, J. Algebra 652 (2024), 52-66. DOI: \href{https://doi.org/10.1016/j.jalgebra.2024.04.008}{10.1016/j.jalgebra.2024.04.008}.  \bibitem{Hreinsdottir2005} F. Hreinsdóttir, \emph{Conjectures on the ring of commuting matrices}, \href{https://arxiv.org/abs/math/0501465}{arXiv:math/0501465}.  \bibitem{Knutson2003} A. Knutson, \emph{Some Schemes Related to the Commuting Variety}, J. Algebraic Geom. 14 (2005), 283-294; preprint \href{https://arxiv.org/abs/math/0306275}{arXiv:math/0306275}.  \bibitem{Youcis2014} A. Youcis, answer to ``What can we say about the image of a regular map?'', Mathematics Stack Exchange, question 731178, answer 732415, 2014. States both the point-set-topology fact that the continuous image of an irreducible space is irreducible, and the \((x,y) \mapsto (x,xy)\) non-closed-image example cited above. CC BY-SA 3.0. \href{https://math.stackexchange.com/questions/731178/what-can-we-say-about-the-image-of-a-regular-map}{math.stackexchange.com/questions/731178}.  \bibitem{Stacks031R} \emph{The Stacks Project}, tag \textbf{031R} = Lemma 10.157.3, Serre's criterion for reducedness: ``Let \(R\) be a Noetherian ring. The following are equivalent: (1) \(R\) is reduced, and (2) \(R\) has properties \(\left( R_{0} \right)\) and \(\left( S_{1} \right)\).'' \href{https://stacks.math.columbia.edu/tag/031R}{stacks.math.columbia.edu/tag/031R}. Note that the enclosing \textbf{section} 10.157 is tag 031O and is titled ``Serre's criterion for \emph{normality}''; its own Lemma 10.157.4 (tag 031S) is the normality criterion \(\left( R_{1} \right)\) and \(\left( S_{2} \right)\), a different statement. 

\bibitem{Tambour2000} T. Tambour, \emph{The Number of Solutions of Some Equations in Finite Groups and a New Proof of Itô's Theorem}, Communications in Algebra \textbf{28} (2000), no. 11, 5353--5361. \href{https://doi.org/10.1080/00927870008827160}{doi:10.1080/00927870008827160}. 

\bibitem{Eisenbud1995} D. Eisenbud, \emph{Commutative Algebra with a View Toward Algebraic Geometry}, Graduate Texts in Mathematics 150, Springer-Verlag, 1995. Cor. 13.4 is the statement used here, equivalent to the Stacks Project tag 00OS above. 

\end{thebibliography}

\end{document}

% Local Variables:
% TeX-engine: xetex
% End:
